% Modernised reading — fonds Alexandre Grothendieck, shelfmark 149, pages 1-135.
%
% This is an INTERPRETATION, not a transcription. It re-reads the pages in
% today's notation and vocabulary, states the hypotheses the manuscript leaves
% implicit, and drops the critical apparatus so that the mathematics reads
% straight through. Everything the brackets and underlines carried moves into a
% footnote. It is held to being mathematically correct as it stands; where the
% manuscript is loose or mistaken the true statement appears and the footnote
% says what the page has. In any dispute of substance the transcription
% governs: whoever cites Grothendieck cites batch-01.fr.tex to batch-07.fr.tex.
%
% Scope: the folder is 135 archivists' pages in seven batches (1-20, ...,
% 121-135), all transcribed, read whole before any of this was written, and
% written as ONE file for the whole shelfmark. Page 1 is a pencilled label
% (« Pages 1-67 ») with no mathematics. His own pagination 1-67 runs on the
% even archivists' pages; the odd pages between continue the text.
%
% The spine. One continuous text in four numbered paragraphs:
%   §1 « Résultats de fidélité » (pages 2-39, his 1-19): the étale topos /
%      fundamental groupoid / Galois group of a field f.g. over Q; Thm 1
%      (trivial centralisers), Scholie, Thm 2 (faithfulness), Cor. 1-3, Thm 3
%      (points give distinct section classes), and a « Complément »;
%   §2 « La question de pleine fidélité » (pages 40-51, his 20-26): the
%      anabelian conjectures (Question-conjecture, Conjecture fondamentale);
%   §3 « Étude des sections de E_U sur Γ » (pages 52-103, his 26-51): points
%      and cusps, Conjectures A, B, B', relative curves, the local criterion,
%      specialisation of sections, bad reduction;
%   §4 « Sections d'extensions, et anneaux de valuation généraux » (pages
%      104-135, his 52-67): notation, inertia, the Riemann-Zariski space, a
%      birational conjecture, the boundary strata. The folder stops mid-§4.
% Reading order: pages 82-83 (his 41) carry formulas (7), (8), which sit
% between (6) of page 66 and (9) of page 72, and continue the argument of
% pages 69-71; they are read there. Page 84 continues page 81.
%
% Notation fixed for the whole folder (his letters in brackets):
%   G_K = Gal(Kbar/K) [E_{Kbar/K}, curly E_K, Gamma when K is algebraic];
%   pi_1(X) the arithmetic fundamental group [E_X, curly E_X]; Delta_X =
%   pi_1(X_Kbar) the geometric one [pi, pi_X, pi_{U/K}]; for a field,
%   Delta_K = Gal(Kbar/K.Qbar) [pi_{Kbar/K}, pi_K]; k the algebraic closure of
%   Q in K, G_k the « arithmetic part » [Gamma_{Kbar/K}, Gamma]; B_K, Pi_K kept;
%   Out for Autext; Zhat(1) = lim mu_n [T, T_infty(Kbar*)]; chi the
%   cyclotomic character and G_K^o = ker chi [Gamma^o, E^o]; L_i the cuspidal
%   inertia (« lacet ») subgroups, as in folder 145; K^h, K^sh fraction fields
%   of the henselisation and strict henselisation [both written K-tilde];
%   « hyperbolic » for his « anabélienne » (2g-2+nu > 0); « polycourbe
%   hyperbolique » for « variété élémentaire (d'Artin) anabélienne ».
%
% Substantive departures from the pages, each footnoted where it occurs:
%   - pp. 12-14: Thm 1 and its corollary on « extensions successives de groupes
%     libres » need the free factors non-abelian (rank >= 2); supplied.
%   - p. 24: 3°) H_1(X) = H_1(J) holds modulo torsion; only surjectivity is
%     used. X normal supplied for the generalised Albanese.
%   - pp. 25, 31: a polycurve of dim >= 2 need NOT embed in a semi-abelian
%     variety (universal elliptic curve minus zero section over Y(N)); the
%     page's claim and its example in Thm 3 hold for curves only.
%   - p. 33: G(K) is not finitely generated when G has a torus part (K* is
%     not); what is true and needed is that G(K) has no non-zero infinitely
%     divisible element.
%   - pp. 36-39: the Lemma, stated without any pi_1 hypothesis, is FALSE for
%     types (1,1) and (1,2) (U = E - {o,-u}; E - {o,t}, t of order 2); the
%     page's argument works for g != 1 and for g = 1, nu >= 3; with the
%     hypothesis of p. 26 all types go through. p. 38: J \ f^{-1}(I), not
%     J \ f(I); the « shrink U » condition does not remove the case card I = 1.
%   - p. 53: the extension (2) is by the NORMALISER of L_i, not its
%     centraliser; the first assertion of the Proposition is false for type
%     (0,2) (every section of G_m is cuspidal) and needs hyperbolicity.
%   - pp. 100-102: c), d) => a) of the « Conjecture de bonne réduction » is
%     false in general (genus 0 with colliding points; genus >= 2 of compact
%     type); a) <=> b) is Oda's criterion.
%   - p. 118: I = Norm_pi(I) can hold only in relative dimension 1.
%   - p. 120: pi_1(U) is G_L modulo the inertia of the divisorial valuations
%     WITHOUT centre on U (the page, partly uncertain, says « ayant un centre »).
%   - p. 121: non-divisorial discrete valuations exist for n >= 2, so codim-0
%     and divisorial points are not the only ones with noetherian local ring.
%   - p. 123: the rank of the inertia is the rational rank of the value group,
%     <= codim (Abhyankar), not = codim.
%   - p. 126: (32)-(33) are read for the model U (section from x in U(K));
%     for E_L as written they contradict pp. 44 and 125.
%   - p. 129: 2°) is false in general (an arc valuation with v(t1) = v(t2)).
%   - p. 133: the boundary map pi_2(D_J) -> T^J is given by Chern classes of
%     normal bundles and can be non-zero ((-1)-curve), which is what the margin
%     of p. 129 senses.
%
% Announced and never established in the folder: the injectivity of
% G_k -> Out(Delta_K) (p. 20, « je sais prouver »); the centraliser fact of
% p. 44 (« à vérifier »); every conjecture of §§2-4; the intersection-
% multiplicity formula of p. 77 (« demander à Carlos » — true, cf. footnote);
% the case nu = 1 of p. 77 (« je vais admettre »); the identification of the
% class c_i with the conormal class (p. 80, sign undecided); the profinite
% characterisation of the geometric splittings (p. 84); e) => a) of p. 102;
% questions 1°), 2°) of p. 129; the struck continuation of p. 135.
%
% Pass: Opus 5.5 (claude-opus-5-5), 2026-10-10 — first pass, unchecked against the pages by a human.
\documentclass[11pt,a4paper]{article}
\input{../preamble/grothendieck}

\folder{149}
\pages{1}{135}
\dating{s.d. — le groupe « Autour de La "Longue Marche" à travers la théorie de Galois » (141 à 149) est daté [à partir de 1978]-1983}
\watermark{Édition de démonstration\\interprétation personnelle de l'œuvre}
\foldertitle{[Autour de La "Longue Marche" à travers la théorie de Galois, pages 1 à 67] : notes manuscrites (s.d.) — lecture modernisée du dossier entier}

\begin{document}

\section*{Résumé}

\begin{resume}
Ces soixante-sept pages, de la numérotation de Grothendieck, forment un texte
suivi, découpé en quatre paragraphes, et posent une question qui paraît
d'abord déraisonnable : peut-on reconstruire une courbe algébrique, ou un
corps de fonctions, définis sur les nombres rationnels, à partir de leurs
seules symétries ?

Les symétries en question sont celles de Galois. À un corps $K$ on associe le
groupe de toutes les permutations de ses extensions algébriques qui respectent
les opérations : son \emph{groupe de Galois absolu}. Quand $K$ contient des
transcendantes — c'est le corps des fonctions d'une courbe, d'une surface —,
ce groupe se coupe en deux. Une partie \emph{géométrique} décrit comment l'on
peut tourner autour des trous de la variété : c'est son groupe fondamental, le
groupe des lacets que le dossier 145 étudie pour la sphère privée de trois
points. Une partie \emph{arithmétique}, le groupe de Galois des nombres
algébriques, agit sur la première. Toute la question du dossier est de savoir
ce que cette action « sait ».

Le premier paragraphe montre qu'elle sait beaucoup : le groupe n'a pas de
centre, si bien que la connaissance de l'action, à automorphisme intérieur
près, suffit ; deux plongements de corps qui induisent la même application
entre groupes sont égaux ; deux points rationnels distincts d'une courbe se
voient comme deux « sections » distinctes du groupe fondamental. La preuve de
ce dernier fait passe par le théorème de Mordell–Weil : la géométrie seule ne
suffit pas, et Grothendieck le dit (« je n'arrive pas à le prouver par voie
géométrique — pour ceci je me rabats sur une méthode arithmétique »). Le
deuxième paragraphe demande la réciproque, « avec un peu de culot » : tout
homomorphisme entre ces groupes qui respecte l'action arithmétique et n'est
pas dégénéré vient-il de la géométrie ? Ce sont les \emph{conjectures
anabéliennes}.

Le troisième paragraphe en isole le cœur, qui est aujourd'hui encore la
pièce manquante : les \emph{sections}. Une image le guide, et c'est la
sienne. Une courbe hyperbolique sur les complexes est le quotient d'un disque
par un groupe de transformations ; son revêtement universel est ce disque,
complété par un cercle de « points à l'infini ». Le groupe de Galois agit sur
ce disque complété. Grothendieck conjecture qu'il y a toujours un point fixe,
et un seul : soit à l'intérieur, et c'est un point rationnel de la courbe ;
soit sur le bord, et c'est une « pointe », un trou rationnel. Un théorème de
point fixe, comme celui de Brouwer pour un disque — mais pour un disque
profini, et l'analogie s'arrête là. Le quatrième paragraphe étend ce
programme aux variétés de dimension quelconque, par les anneaux de valuation
et l'espace de Riemann–Zariski ; il s'interrompt au milieu d'une construction.

Le texte dit aussi, en passant, comment il travaille. Il essaie, échoue, se
rabat (« Ouf ! » à la fin d'une démonstration de quatre pages) ; il hasarde
des conjectures qu'il qualifie lui-même d'« hérétiques », puis les retrouve
dans la logique de ce qui suit ; et il s'impatiente du langage même dans
lequel il écrit : il faudra, dit-il, « développer un langage géométrique qui
colle mieux à l'intuition géométrique, que les sempiternelles extensions de
groupes profinis ».

Ces conjectures sont celles que Grothendieck énoncera dans sa lettre à
Faltings de juin 1983. Une partie en a été démontrée dans les années 1990
(Nakamura, Tamagawa, Pop, Mochizuki) : on reconstruit effectivement une
courbe hyperbolique, ou un corps de fonctions, à partir de son groupe
fondamental muni de l'action de Galois. La \emph{conjecture des sections},
elle, reste ouverte à notre connaissance. Pour aller plus loin : géométrie
anabélienne, conjecture des sections, sections cuspidales et points base
tangentiels, polycourbes hyperboliques, critère de bonne réduction d'Oda,
espace de Riemann–Zariski, complexe dual d'un diviseur à croisements normaux.

\keywords{étale fundamental group, absolute Galois group, outer Galois representation, anabelian geometry, Grothendieck anabelian conjecture, Hom conjecture, faithfulness, centraliser, cohomological dimension, Artin neighbourhood, hyperbolic polycurve, generalised Jacobian, Kummer theory, Mordell–Weil theorem, Lefschetz fixed-point formula, section conjecture, cuspidal section, tangential base point, birational section conjecture, universal cover, specialisation of sections, purity, decomposition group, inertia group, henselisation, good reduction criterion, Zariski–Riemann space, valuation ring, Abhyankar inequality, normal crossings divisor, dual complex}
\end{resume}

\pagerange{2}{135}
\section*{Le fil du dossier, et les conventions}

\subsection*{Les stations}

\begin{itemize}
\item \textbf{§1, « Résultats de fidélité » (pages 2 à 39, ses pages 1 à 19).}
Le topos étale, le groupoïde fondamental et le groupe de Galois d'un corps de
type fini sur $\mathbb{Q}$ ; partie géométrique et partie arithmétique
(pages 2 à 11) ; dimension cohomologique (11-12) ; théorème 1, centralisateurs
triviaux (12-14) ; scholie et descriptions équivalentes (14-19) ; un
renforcement (20) ; théorème 2, fidélité, avec les corollaires 1 à 3 et le
théorème 3 (21-35) ; un « Complément » géométrique (36-39).
\item \textbf{§2, « La question de pleine fidélité » (pages 40 à 51, ses pages
20 à 26).} Les homomorphismes « admissibles » ; la question-conjecture sur les
corps (40-44) ; les modèles et la « conjecture fondamentale » (45-49) ; un
programme de reconstruction (49-51).
\item \textbf{§3, « Étude des sections de $\pi_1(U)$ sur $G_K$ » (pages 52 à
103, ses pages 26 à 51).} Points et pointes, conjecture A (52-57) ; la forme
« point fixe », conjectures B et B$'$ (57-62) ; courbes relatives et
reconstruction (63-68) ; sections rationnelles et sections partout définies
(69-78, avec le feuillet 82-83) ; la classe de l'extension de seconde espèce
(78-84) ; spécialisation des sections (84-98) ; mauvaise réduction (98-103).
\item \textbf{§4, « Sections d'extensions, et anneaux de valuation généraux »
(pages 104 à 135, ses pages 52 à 67).} Révision des notations (104-114) ;
modèles et inertie (114-120) ; l'espace de Riemann–Zariski (120-123) ; germes
de sections, une conjecture birationnelle (123-127) ; l'inertie le long du bord
(127-131) ; strates du bord (131-135). Le dossier s'arrête là, sur un
paragraphe barré.
\end{itemize}

\emph{L'ordre des feuillets.} Grothendieck numérote une page sur deux (ses
pages 1 à 67 sont les pages paires 2 à 134 des archivistes) ; les pages
impaires continuent le texte sans rupture. Une seule exception : les pages 82
et 83 (sa page 41) ne prolongent pas la page 81 ; elles portent les formules
numérotées (7) et (8), qui s'insèrent entre la formule (6) de la page 66 et la
formule (9) de la page 72, et elles poursuivent l'argument des pages 69 à 71.
On les lit là. La page 84 reprend exactement où s'arrête la page 81. Cette
remise en ordre est la nôtre ; la numérotation des formules est ce qui la
fonde.

\emph{Ce qui n'est pas établi.} Le dossier démontre les énoncés du §1, à
quelques points près que les notes signalent. Tout ce qui suit est
conjectural, et Grothendieck le dit à chaque pas ; quelques présomptions des
§§3 et 4 sont fausses telles qu'écrites, et les notes le montrent par un
exemple.

\subsection*{Les conventions}

\emph{Corps et groupes de Galois.} $K$ est un corps de type fini sur
$\mathbb{Q}$, de degré de transcendance $n$ ; $\overline{K}$ une clôture
algébrique ; $G_K = \mathrm{Gal}(\overline{K}/K)$\footnote{Il écrit
$E_{\overline{K}/K}$, puis $\mathcal{E}_K$ (une lettre ronde), et $\Gamma$
quand le corps est algébrique sur $\mathbb{Q}$ ; au §4, $\mathcal{E}^{\overline
K}_K$ et $\Gamma_K$. Nous écrivons $G_K$ partout.}. On note $k$ la fermeture
algébrique de $\mathbb{Q}$ dans $K$, extension finie de $\mathbb{Q}$, et
$\overline{\mathbb{Q}}$ celle de $\mathbb{Q}$ dans $\overline{K}$. Le
caractère cyclotomique est $\chi : G_K \to \widehat{\mathbb{Z}}^{*}$, et
$G_K^{\circ} = \ker \chi = G_{K(\mu_\infty)}$\footnote{Il écrit
$\Gamma^{\circ}$, $E^{\circ}$, $\mathcal{E}_K^{\circ}$.}. On écrit
$\widehat{\mathbb{Z}}(1) = \varprojlim_n \mu_n$, module de Tate du groupe
multiplicatif\footnote{Il écrit $T$, ou $T_\infty(\overline{K}^{*})$, ou
$\varprojlim \mu_n$.}.

\emph{Groupes fondamentaux.} Pour un $K$-schéma $X$ géométriquement connexe,
$\pi_1(X)$ est son groupe fondamental étale (arithmétique), et
$\Delta_X = \pi_1(X_{\overline{K}})$ sa partie géométrique, d'où la suite
exacte fondamentale
\[
1 \longrightarrow \Delta_X \longrightarrow \pi_1(X) \longrightarrow G_K
\longrightarrow 1 .
\]
Pour $X = \operatorname{Spec} K$ on écrit $\Delta_K = \mathrm{Gal}(\overline{K}
/ K\overline{\mathbb{Q}})$\footnote{Il écrit $\mathcal{E}_X$ ou $E_X$ pour
$\pi_1(X)$, et $\pi$, $\pi_X$, $\pi_{U/K}$, $\pi_{\overline{K}/K}$, $\pi_K$
pour la partie géométrique.}. Les groupes « extérieurs » sont les groupes
définis à automorphisme intérieur près ; $\mathrm{Out}(\Delta) =
\mathrm{Aut}(\Delta)/\mathrm{Int}(\Delta)$\footnote{Son $\mathrm{Autext}$. Le
dossier 141 (pages 11 à 15) développe la même notion sous forme de
bitorseurs, $\mathrm{Isomext}(\pi', \pi)$.}. $B_K$ est le topos étale de
$\operatorname{Spec} K$, $\Pi_K$ son groupoïde fondamental : ses objets sont
les clôtures séparables de $K$.

\emph{Courbes.} Une courbe de type $(g, \nu)$ est une courbe lisse
géométriquement connexe, complémentaire de $\nu$ points dans une courbe
propre $X = \widehat{U}$ de genre $g$ ; elle est \emph{hyperbolique} si
$2g - 2 + \nu > 0$\footnote{Son mot est « anabélienne » : $\Delta_U$ est
alors non abélien, et de centre trivial. Il exclut ainsi les types $(0,0)$,
$(0,1)$, $(0,2)$, $(1,0)$.}. Les points de $X \setminus U$ sont ses
\emph{pointes} ; à chaque pointe $i$ et à chaque point de
$\widetilde{U}$ au-dessus correspond un sous-groupe d'inertie $L_i \subset
\Delta_U$, isomorphe à $\widehat{\mathbb{Z}}(1)$, bien défini à conjugaison
près — le « sous-groupe de lacets » de Grothendieck, noté $L_i$ comme dans
le dossier 145.

\emph{Variétés élémentaires.} Une \emph{variété élémentaire} (d'Artin) est un
schéma lisse qui se dévisse en fibrations successives de courbes lisses
(fibrations élémentaires de SGA 4, exposé XI) ; elle est \emph{hyperbolique}
si toutes les fibres le sont — on dit aujourd'hui \emph{polycourbe
hyperbolique}\footnote{Son « variété élémentaire anabélienne », ou
« admissible ». Le nom de polycourbe est postérieur (Hoshi, Schmidt–Stix).}.

\emph{Sections.} Une section de $\pi_1(U) \to G_K$ est un homomorphisme
continu $s : G_K \to \pi_1(U)$ relevant l'identité ; on la considère à
conjugaison par $\Delta_U$ près. Un point $x \in U(K)$ en donne une classe :
ce sont les sections \emph{de première espèce}. Une pointe rationnelle en
donne d'autres, par le normalisateur de $L_i$ : ce sont les sections
\emph{de seconde espèce}, aujourd'hui \emph{sections cuspidales}\footnote{Il
emploie lui-même « de 2\textsuperscript{e} espèce » à la page 78.}.

\emph{Hensélisés.} Pour un point $s$ d'un schéma normal, $K^{h}_s$ et
$K^{sh}_s$ sont les corps des fractions de l'hensélisé et de l'hensélisé
strict de l'anneau local\footnote{Il écrit $\tilde K$ pour les deux, à une
ligne d'intervalle (page 92).}.

\pagerange{2}{39}
\section*{§1. Résultats de fidélité (pages 2 à 39)}

\pagerange{2}{11}
\subsection*{Le groupe de Galois d'un corps, partie géométrique et partie arithmétique (pages 2 à 11)}

À tout corps $K$ on associe son topos étale $B_K$, topos galoisien
profini, et le groupoïde $\Pi_K$ de ses points : ce groupoïde est
canoniquement anti-équivalent à celui des clôtures séparables de $K$. Pour une
telle clôture $\overline{K}$, le groupe d'automorphismes du point
correspondant s'identifie à $G_K$, ou plutôt à son opposé\footnote{La
parenthèse de la page (« il vaut peut-être mieux dire : à l'opposé de ce
groupe ») tient à ce que les clôtures varient comme les foncteurs fibres, à
l'inverse des points. Le choix n'a pas d'incidence sur ce qui suit, puisque
$g \mapsto g^{-1}$ identifie un groupe à son opposé.}. On retrouve $B_K$
soit comme le topos des systèmes locaux sur $\Pi_K$, soit comme celui des
ensembles discrets munis d'une action continue de $G_K$.

Un plongement $K \to K'$ définit un morphisme de topos $B_{K'} \to B_K$ et
un morphisme de groupoïdes $\Pi_{K'} \to \Pi_K$ : une clôture $\overline{K'}$
de $K'$ en définit une de $K$, la fermeture algébrique de $K$ dans
$\overline{K'}$, et l'on obtient un homomorphisme
\[
G_{K'} \longrightarrow G_K
\]
dont l'image est le sous-groupe fermé qui fixe $K_1 = \overline{K} \cap K'$,
fermeture algébrique de $K$ dans $K'$\footnote{La page écrit $K^{*}$ pour
$K$ en trois endroits ; la petite marque ne paraît désigner que $K$
lui-même.}. Si $K'$ est de type fini sur $K$, $K_1$ est fini sur $K$ :
l'image est ouverte, et égale à $G_K$ si et seulement si $K$ est
séparablement clos dans $K'$. L'homomorphisme est injectif si et seulement si
$K'$ est algébrique sur $K$, donc bijectif si et seulement si $K'/K$ est
radicielle — en caractéristique $0$, si et seulement si $K \to K'$ est un
isomorphisme. Ainsi les foncteurs $K \mapsto B_K$, $K \mapsto \Pi_K$,
$K \mapsto G_K$ sont \emph{conservatifs} sur la catégorie des corps de
caractéristique $0$ et des extensions de type fini\footnote{Une note en
diagonale, en partie illisible, propose pour la caractéristique $p$ de
passer à la catégorie quotient où les extensions radicielles finies sont
inversées ; c'est la forme juste de l'énoncé en caractéristique $p$.}.

Pour des corps quelconques, ces foncteurs sont loin d'être fidèles : tous les
corps séparablement clos ont le même topos, le topos ponctuel. On associe donc
à $K$ un objet plus fin, le système projectif $\mathbb{B}_K = (B_{K_i})$, où
$K_i$ parcourt les sous-corps de $K$ de type fini sur le corps premier, et de
même $\boldsymbol{\Pi}_K$ et $\mathbb{G}_K$. On a $B_K \simeq \varprojlim
B_{K_i}$, $\Pi_K \simeq \varprojlim \Pi_{K_i}$, $G_K \simeq \varprojlim
G_{K_i}$ ; la réciproque ne vaut pas, mais comme le foncteur
$\mathrm{Ind}(\text{corps de type fini}) \to \text{corps}$ est une
équivalence, les foncteurs pro-objets $K \mapsto \mathbb{B}_K$, etc., ont
exactement les propriétés de fidélité de leurs restrictions aux corps de type
fini, auxquels on se borne désormais. Il faudra pourtant, en cours de route,
une description algébrique des pro-groupes attachés par exemple à
$\mathbb{C}$.

Le rôle dominant est tenu par $\mathbb{Q}$, et par $G_{\mathbb{Q}} =
\mathrm{Gal}(\overline{\mathbb{Q}}_0/\mathbb{Q})$, où
$\overline{\mathbb{Q}}_0$ est la clôture algébrique de $\mathbb{Q}$ dans
$\mathbb{C}$. Tout corps $K$ de caractéristique $0$ est muni de $B_K \to
B_{\mathbb{Q}}$, $\Pi_K \to \Pi_{\mathbb{Q}}$, c'est-à-dire, une clôture
$\overline{K}$ étant choisie, d'un homomorphisme $G_K \to
G_{\mathbb{Q}}$, et l'on regarde désormais $B_K$, $\Pi_K$, $G_K$ comme munis
de cette structure. Son noyau est la \emph{partie géométrique}
$\Delta_K$ ; son image, la \emph{partie arithmétique}, est le sous-groupe
ouvert $G_k \subset G_{\mathbb{Q}}$, où $k$ est la fermeture algébrique de
$\mathbb{Q}$ dans $K$ :
\[
1 \longrightarrow \Delta_K \longrightarrow G_K \longrightarrow G_k
\longrightarrow 1 .
\]
On a $\Delta_K = 1$ si et seulement si $K$ est fini sur $\mathbb{Q}$.

Pour interpréter $\Delta_K$, on écrit $K = \varinjlim A_i$, où $A_i$
parcourt les sous-$\mathbb{Q}$-algèbres de type fini de $K$ ; les $U_i =
\operatorname{Spec} A_i$ sont les \emph{modèles affines} de $K$. On peut se
borner à des $A_i$ réguliers, reliés par des localisations $A_j = (A_i)_{f}$,
donc à des $U_i$ lisses sur $k$ et géométriquement intègres, reliés par des
immersions ouvertes ; et même, d'après Artin, à des variétés élémentaires.
Le choix de $\overline{K}$ définit un point géométrique $\bar\eta$ commun aux
$U_i$, et l'on a
\[
G_K = \pi_1(\operatorname{Spec} K, \bar\eta) \xrightarrow{\ \sim\ }
\varprojlim_i \pi_1(U_i, \bar\eta), \qquad
\Delta_K \xrightarrow{\ \sim\ } \varprojlim_i \Delta_{U_i}, \qquad
\Delta_{U_i} = \pi_1(U_i \otimes_k \overline{\mathbb{Q}}, \bar\eta),
\]
la seconde égalité parce que la suite exacte de $K$ est la limite projective
des suites exactes d'homotopie
\[
1 \longrightarrow \Delta_{U_i} \longrightarrow \pi_1(U_i) \longrightarrow
G_k \longrightarrow 1 .
\]
Par un plongement $\overline{K} \subset \mathbb{C}$, $\Delta_{U_i}$ est le
complété profini du groupe fondamental de la variété complexe
$U_i(\mathbb{C})$ (théorème d'existence de Riemann)\footnote{La page note
$\Gamma_i$ dans une parenthèse ce qu'elle appelle $E_i$ dans le texte ; un
membre de l'égalité placée sous $\pi_1(k, \bar\eta)$ est surchargé et
illisible ; une note en diagonale, très pâle, renvoie aux formules (12) et (13)
et à « $K/k_0$ » sans qu'on puisse en lire davantage.}.

\pagerange{11}{12}
\subsection*{Dimension cohomologique (pages 11 et 12)}

Si $U_i$ est une variété élémentaire de dimension $n$, les fibrations
successives font de $\Delta_{U_i}$ une extension successive de $n$ groupes
profinis libres (ceux des fibres, courbes affines). D'où $\mathrm{cd}\,
\Delta_{U_i} = n$, $\mathrm{cd}_\ell\, \pi_1(U_i) \le n+2$ pour $\ell$ impair,
et à la limite
\[
\mathrm{cd}\, \Delta_K \le n, \qquad \mathrm{cd}_\ell\, G_K \le n + 2 \quad
(\ell \ne 2),
\]
qui sont en fait des égalités\footnote{La page ajoute « s.f. erreur ». Les
égalités sont exactes : $\Delta_K$ est le groupe de Galois absolu de
$K\overline{\mathbb{Q}}$, de type fini et de degré de transcendance $n$ sur
un corps algébriquement clos, de dimension cohomologique $n$ (Serre,
\emph{Cohomologie galoisienne}, II.4.2) ; pour $\ell$ impair,
$\mathrm{cd}_\ell\, G_k = 2$, et la formule d'additivité donne $n+2$. Le
numéro de la formule, lu « (14) », ressemble sur la page à « (K) ».} : le
degré de transcendance absolu de $K$ se lit sur la cohomologie de $G_K$.

\pagerange{12}{14}
\subsection*{Théorème 1 : centralisateurs triviaux (pages 12 à 14)}

\textbf{Théorème 1.} \emph{Soit $K$ de type fini sur $\mathbb{Q}$. Le
centralisateur dans $G_K$ de tout sous-groupe ouvert de $G_K$ est trivial ;
de même dans $\Delta_K$.}

La démonstration passe par deux corollaires, qui en sont les deux cas
extrêmes.

\textbf{Corollaire} (le cas arithmétique)\textbf{.} \emph{Dans
$G_{\mathbb{Q}}$, le centralisateur de tout sous-groupe ouvert est
trivial.} Soit $\Gamma'$ ouvert, qu'on peut supposer distingué, et $Z$ son
centralisateur. Le centre de $\Gamma'$ est trivial\footnote{La page dit « il
est bien connu (?) » et ajoute en marge « ? à vérifier ». C'est vrai : le
groupe de Galois absolu d'un corps de nombres est de centre trivial.}, donc
$Z \cap \Gamma' = 1$, $Z$ s'injecte dans $G_{\mathbb{Q}}/\Gamma'$ et est fini.
Or les seuls éléments d'ordre fini non triviaux de $G_{\mathbb{Q}}$ sont les
conjugués de la conjugaison complexe $\tau$ (Artin–Schreier), et le
centralisateur de $\tau$ est $\{1, \tau\}$, qui ne contient pas $\Gamma'$. Donc
$Z = 1$.

\textbf{Corollaire} (le cas géométrique)\textbf{.} \emph{Soit $\Delta$ un
groupe profini, extension successive de groupes profinis libres non
abéliens. Le centralisateur dans $\Delta$ de tout sous-groupe ouvert est
trivial.}\footnote{La page dit « libres », sans « non abéliens » ;
l'hypothèse est indispensable, puisque $\widehat{\mathbb{Z}}$, libre de rang
$1$, est abélien. Elle est satisfaite en choisissant des modèles dont les
fibres sont hyperboliques, ce qu'on peut toujours faire en ôtant des points :
c'est l'hypothèse « anabélienne » que la page ajoute plus loin (pages 23 et
27).} Par récurrence : si $N \subset \Delta$ est le premier facteur libre et
$Z$ centralise $\Delta'$ ouvert, l'image de $Z$ dans $\Delta/N$ centralise
l'image, ouverte, de $\Delta'$, donc est triviale ; ainsi $Z \subset N$ et $Z$
centralise $\Delta' \cap N$, ouvert dans $N$. On est ramené à un groupe libre
$F$ de rang $\ge 2$ et à $\Delta'$ ouvert distingué : $\Delta'$ est libre de
rang $\ge 2$, donc de centre trivial, $Z$ s'injecte dans $F/\Delta'$ et est
fini, et un groupe profini libre est sans torsion\footnote{La page admet ces
deux faits, en marge : « à vérifier ». Ils sont classiques ; le second tient à
ce qu'un groupe profini libre est de dimension cohomologique $1$.}. Le cas
d'une limite projective $\Delta_K = \varprojlim \Delta_{U_i}$, à applications
de transition surjectives, s'en déduit aussitôt.

\emph{Démonstration du théorème.} Si $Z$ centralise $E' \subset G_K$ ouvert,
son image dans $G_k$ centralise l'image de $E'$, ouverte, donc est triviale
par le premier corollaire : $Z \subset \Delta_K$, et $Z$ centralise $E' \cap
\Delta_K$, ouvert dans $\Delta_K$, donc $Z = 1$ par le second\footnote{Le
début de la démonstration, page 12, est enchevêtré de trois lignes biffées ;
seul le sens général se lit, celui qu'on donne.}.

\pagerange{14}{19}
\subsection*{Scholie : ce que voit le groupe extérieur (pages 14 à 19)}

Parce que $G_K$ est de centre trivial, le groupoïde $\Pi_K$ — et donc
$B_K$ — est connu à isomorphisme unique près dès qu'on connaît le groupe
extérieur sous-jacent à $G_K$. Les homomorphismes $G_{K'} \to G_K$ associés
aux extensions de type fini ont une image ouverte, donc de centralisateur
trivial ; par suite les morphismes $B_{K'} \to B_K$ et $\Pi_{K'} \to \Pi_K$
sont déterminés, à isomorphisme unique près, par l'homomorphisme extérieur
correspondant. Il en va ainsi des morphismes structuraux $B_K \to
B_{\mathbb{Q}}$.

Grothendieck préfère un point de vue voisin, qui exploite le fait que
$\Delta_K$ est lui aussi de centre trivial : l'extension de $G_k$ par
$\Delta_K$ est alors entièrement déterminée, à isomorphisme unique près, par
l'action extérieure
\[
\varphi : G_k \longrightarrow \mathrm{Out}(\Delta_K),
\]
comme image inverse par $\varphi$ de l'extension universelle

\begin{tikzcd}
1 \arrow[r] & \Delta_K \arrow[r] & \mathrm{Aut}(\Delta_K) \arrow[r] & \mathrm{Out}(\Delta_K) \arrow[r] & 1 \\
 & & & G_k \arrow[u, "\varphi"'] &
\end{tikzcd}

Un couple $(K, \overline{K})$ est donc décrit, avec une fidélité qui reste à
examiner, par un triple $(\Delta, \Gamma, \varphi)$ — groupe profini $\Delta$,
sous-groupe ouvert $\Gamma$ d'un $G_{\mathbb{Q}}$ bien déterminé, action
extérieure $\varphi$ — d'où l'on reconstruit l'extension $G_K$ et
l'homomorphisme $\Pi_K \to \Pi_{\mathbb{Q}}$ avec un objet de
$\Pi_K$\footnote{« J'ai oublié de préciser », écrit-il, qu'il faut se donner
$\Gamma$ comme sous-groupe d'un $G_{\mathbb{Q}}$ déterminé, c'est-à-dire un
objet de $\Pi_{\mathbb{Q}}$ sur lequel $\Gamma$ opère fidèlement.}. On
obtient ainsi plusieurs descriptions essentiellement équivalentes d'une
extension $K$ de type fini de $\mathbb{Q}$ :
\begin{enumerate}
\item le topos $B_K$, comme topos pro-galoisien sur $B_{\mathbb{Q}}$ ;
\item son groupoïde fondamental $\Pi_K$, au-dessus de $\Pi_{\mathbb{Q}}$ ;
\item le groupe extérieur $G_K$, au-dessus du groupe extérieur
$G_{\mathbb{Q}}$ ;
\item pour un couple $(K, \overline{K})$ : un objet $\overline{\mathbb{Q}}$ de
$\Pi_{\mathbb{Q}}$ et un homomorphisme $G_K \to
\mathrm{Gal}(\overline{\mathbb{Q}}/\mathbb{Q})$ ;
\item une clôture $\overline{\mathbb{Q}}$ étant fixée, pour un couple $(K,
i)$ où $i : k \to \overline{\mathbb{Q}}$ : le groupe extérieur $\Delta_K$ sur
lequel un sous-groupe ouvert de $G_{\mathbb{Q}}$ opère extérieurement ;
\item (en marge) le groupoïde de $K \otimes_{\mathbb{Q}} \overline{\mathbb{Q}}$,
ou son topos galoisien, sur lequel $G_{\mathbb{Q}}$ opère\footnote{Note
verticale en marge, en partie illisible. La page 19 écrit aussi « $\pi_1(K) =
\Gamma_K$ » là où l'on attend $\Delta_K$.}.
\end{enumerate}
Dans la description 5, un plongement $K \to K'$ donne un homomorphisme
extérieur $\Delta_{K'} \to \Delta_K$ d'image ouverte, donc de centralisateur
trivial, commutant aux actions extérieures de $G_{k'}$ ; ces données suffisent
à reconstituer les extensions $G_K$, $G_{K'}$, les topos à équivalence près,
et l'homomorphisme d'extensions.

\pagerange{20}{20}
\subsection*{Remarque : un renforcement (page 20)}

Si $K$ n'est pas fini sur $\mathbb{Q}$, le théorème 1 se renforce : pour
tout $\Delta' \subset \Delta_K$ ouvert \emph{dans $\Delta_K$}, on a
$\mathrm{Centr}_{G_K}(\Delta') = 1$, pourvu que l'action extérieure $\varphi :
G_k \to \mathrm{Out}(\Delta_K)$ soit injective. En effet, si $Z$ est ce
centralisateur, $Z \cap \Delta_K = 1$ par le théorème 1 ; quitte à
remplacer $K$ par une extension finie, on peut supposer $\Delta' = \Delta_K$,
et alors l'image de $Z$ dans $G_k$ agit trivialement sur $\Delta_K$, à
automorphisme intérieur près, donc est dans $\ker \varphi = 1$, et $Z = 1$.
Grothendieck écrit : « sauf erreur je sais prouver que cet homomorphisme est
injectif (on est ramené aussitôt au cas où $K$ est de degré de transcendance
1, et au cas du $\pi_1$ d'une courbe algébrique…) », et la page s'arrête
là\footnote{La page annonçait d'abord une conclusion plus faible : $Z$
d'ordre $1$ ou $2$, engendré dans le second cas par une conjugaison
complexe. L'injectivité de la représentation galoisienne extérieure d'une
courbe hyperbolique sur un corps de nombres est aujourd'hui établie : pour
$\mathbb{P}^1 \setminus \{0, 1, \infty\}$ c'est une conséquence du
théorème de Belyi (1979) ; pour les courbes hyperboliques affines, puis
propres, elle a été démontrée dans les années 1990 et 2000 (Matsumoto,
Hoshi–Mochizuki). Rien ne dit quel argument il avait en vue.}.

\pagerange{21}{39}
\section*{Théorème 2 : fidélité (pages 21 à 39)}

\pagerange{21}{23}
\subsection*{L'énoncé, et sa réduction à la géométrie (pages 21 à 23)}

\textbf{Théorème 2.} \emph{Le foncteur $K \mapsto (\Pi_K \to \Pi_{\mathbb{Q}})$,
des extensions de type fini de $\mathbb{Q}$ vers les groupoïdes profinis
au-dessus de $\Pi_{\mathbb{Q}}$, est fidèle.}

Autrement dit, si deux plongements $f, g : K \rightrightarrows K'$ induisent
des morphismes isomorphes au-dessus de $\Pi_{\mathbb{Q}}$, alors $f = g$. En
termes de groupes : une clôture $\overline{K'}$ étant choisie, $f$ et $g$ en
déduisent deux clôtures de $K$, et l'hypothèse dit qu'on peut les identifier
de sorte que $f^{*} = g^{*} : G_{K'} \to G_K$ ; une clôture
$\overline{\mathbb{Q}}$ étant fixée, que les deux homomorphismes extérieurs
$\Delta_{K'} \rightrightarrows \Delta_K$, compatibles à l'action de
$G_{\mathbb{Q}}$, sont égaux\footnote{Une note en marge observe que le
théorème « n'est pas spécial à $\mathbb{Q}$ » et vaudrait sur un corps de
base quelconque ; la suite de la note est illisible. Une autre, en bas de
marge, ramène au cas $K = K'$.}.

\emph{Réduction.} Écrivons $\operatorname{Spec} K = \varprojlim U_i$, d'où
$\Delta_K = \varprojlim \Delta_{U_i}$. Il suffit de voir que $f_i = f|_{A_i}$
et $g_i = g|_{A_i}$ coïncident pour tout $i$. Pour $i$ fixé, $K' =
\varinjlim A'_j$ avec $A'_j \supset f(A_i) \cup g(A_i)$, et $\Delta_{K'} =
\varprojlim \Delta_{V_j}$, $V_j = \operatorname{Spec} A'_j$. Les $V_j$ étant
réguliers, les applications $\Delta_{K'} \to \Delta_{V_j}$ sont surjectives :
l'égalité des homomorphismes extérieurs $\Delta_{K'} \to \Delta_{U_i}$
entraîne donc celle de $\Delta_{V_j} \to \Delta_{U_i}$. Les morphismes $f_i,
g_i : V_j \rightrightarrows U_i$ sont dominants, puisque $f$ et $g$ sont des
plongements de corps ; et l'on est ramené à un énoncé « purement
géométrique », sur un corps algébriquement clos.

\textbf{Corollaire 1.} \emph{Soient $X$, $Y$ des schémas de type fini sur un
corps algébriquement clos $k$, $X$ normal et connexe, et $f, g : X
\rightrightarrows Y$ tels que $\pi_1(f) = \pi_1(g)$ comme homomorphismes
extérieurs.}
\begin{enumerate}
\item[a)] \emph{Si $Y$ se plonge, par $i$, dans un groupe algébrique
commutatif $G$, extension d'une variété abélienne par un tore, il existe un
unique $a \in G(k)$ tel que $i \circ g = \tau_a \circ i \circ f$, où $\tau_a$
est la translation par $a$.}
\item[b)] \emph{Si $Y$ est une polycourbe hyperbolique et si $f$ ou $g$ est
dominant, alors $f = g$.}
\end{enumerate}
\footnote{La page dit d'abord « $X$, $Y$ réduits », puis, pour b), « $X$
réduit » sur un « normal » biffé. Nous supposons $X$ normal et connexe, ce
qu'exige la jacobienne généralisée utilisée en a), et ce que le corollaire 2
supposera. Une note marginale oblique, rattachée à b), est en grande partie
illisible. Une autre observe que dans a) il suffit que $H_1(f, \mathbb{Z}_\ell)
= H_1(g, \mathbb{Z}_\ell)$ pour un $\ell$ premier à la caractéristique.}

\pagerange{23}{24}
\subsection*{Le cas abélien (pages 23 et 24)}

L'unicité de $a$ est claire : $a = g(x) - f(x)$ pour tout $x$. Pour
l'existence, on remplace $Y$ par $G$, on choisit une origine $x_0 \in X(k)$ et
l'on se ramène à $f(x_0) = g(x_0) = 0$ ; il s'agit alors de prouver $f = g$, et
l'on peut affaiblir l'hypothèse en $H_1(f, \mathbb{Z}_\ell) = H_1(g,
\mathbb{Z}_\ell)$ pour un $\ell$ premier à la caractéristique. Soit $J$ la
jacobienne généralisée de $X$ (sa variété d'Albanese généralisée, extension
d'une variété abélienne par un tore). On sait que
\begin{enumerate}
\item[1°)] tout morphisme $f : X \to G$ tel que $f(x_0) = 0$ se factorise de
façon unique en $X \xrightarrow{\mathrm{can}} J \xrightarrow{\varphi} G$, avec
$\varphi$ un homomorphisme de groupes algébriques ;
\item[2°)] un tel $\varphi$ est déterminé par son action sur $H_1(-,
\mathbb{Z}_\ell) = T_\ell$, c'est-à-dire sur les points de torsion
$\ell$-primaire, qui sont denses ;
\item[3°)] $H_1(X, \mathbb{Z}_\ell) \to H_1(J, \mathbb{Z}_\ell)$ est
surjectif\footnote{La page écrit un isomorphisme. C'en est un modulo torsion
(en caractéristique $0$) ; seule la surjectivité sert.}.
\end{enumerate}
Si $f = \varphi \circ \mathrm{can}$ et $g = \psi \circ \mathrm{can}$, l'égalité
$H_1(f) = H_1(g)$ et 3°) donnent $H_1(\varphi) = H_1(\psi)$, donc $\varphi =
\psi$ par 2°), et $f = g$.

L'hypothèse de a) ne peut donner $f = g$ sans normalisation : pour $X = Y = G$,
les translations induisent l'identité sur $\pi_1(G)$ sans être l'identité.

\pagerange{24}{27}
\subsection*{Le cas hyperbolique : réduction à un énoncé sur les translations (pages 24 à 27)}

L'hypothèse de dominance n'est pas superflue : deux applications constantes
distinctes induisent la même application (triviale) sur les $\pi_1$.

Supposons d'abord $Y$ plongé par $i$ dans une extension $G$ d'une variété
abélienne par un tore\footnote{La page affirme qu'une variété élémentaire
« se plonge bien sûr » dans un tel groupe. C'est vrai pour une courbe
hyperbolique (dans sa jacobienne généralisée), faux en général en dimension
$\ge 2$ : la courbe elliptique universelle privée de sa section nulle,
au-dessus d'une courbe modulaire $Y(N)$, $N \ge 3$, est une polycourbe
hyperbolique de dimension $2$ dont la monodromie sur le $H_1$ de la fibre n'a
que des coinvariants finis ; son application d'Albanese généralisée se
factorise par la base et n'est pas un plongement. Les arguments des pages 25 à
39 valent donc tels quels pour les courbes ; en dimension supérieure ils
demandent un dévissage le long de la fibration, que la page invoque page 36
(« par dévissage ») sans l'écrire.}. Par a), $i \circ g = \tau_u \circ i \circ
f$ pour un $u \in G(k)$. Comme $f$ est dominant, il existe un ouvert non vide
$V \subset Y$ contenu dans l'image de $f$ et tel que $V + u \subset Y$ ; quitte
à remplacer $X$ par un ouvert dense, $f$ se factorise par $j : V \hookrightarrow
Y$, et $g$ par $j' : V \to Y$, $y \mapsto y + u$. L'image de $\pi_1(X) \to
\pi_1(V)$ est un sous-groupe \emph{ouvert} (c'est vrai pour un morphisme
dominant entre variétés normales). L'hypothèse dit donc que $\pi_1(j)$ et
$\pi_1(j')$ coïncident, à automorphisme intérieur de $\pi_1(Y)$ près, sur un
sous-groupe ouvert de $\pi_1(V)$, et l'on est ramené à :

\begin{quote}
\emph{($\ast$) Soient $V \subset Y$ un ouvert non vide et $u \in G$ tels que
$V + u \subset Y$, et que $\pi_1(j)$, $\pi_1(j')$ coïncident extérieurement sur
un sous-groupe ouvert de $\pi_1(V)$. Alors $u = 0$.}
\end{quote}

\footnote{Dans ces lignes la page note $U$ l'ouvert de $Y$, d'abord noté $V$,
comme l'ouvert de $X$ ; nous gardons $V$. Une note en marge observe que ($\ast$)
n'est pas plus fort que le corollaire : on le retrouve en prenant pour $X$ le
revêtement étale fini de $V$ attaché au sous-groupe ouvert.}« Finalement, je
n'arrive pas à le prouver par voie \emph{géométrique} — pour ceci je me
rabats sur une méthode \emph{arithmétique} » : ce sont les corollaires 2 et 3
et le théorème 3. La démonstration géométrique viendra quand même, dans le
« Complément » des pages 36 à 39.

\pagerange{27}{33}
\subsection*{Le détour arithmétique : corollaires 2 et 3, théorème 3 (pages 27 à 33)}

\textbf{Corollaire 2.} \emph{La conclusion $f = g$ de 1 b) reste vraie si,
au lieu de supposer $f$ ou $g$ dominant, on suppose $k$ de caractéristique
$0$, $X$ normal, et l'une des hypothèses suivantes :}
\begin{enumerate}
\item[c)] \emph{l'image de $\pi_1(f)$ est ouverte dans $\pi_1(Y)$, et $Y$ est
une polycourbe hyperbolique ;}
\item[d)] \emph{l'image de $\pi_1(f)$ a un centralisateur trivial dans
$\pi_1(Y)$, et $Y$ se plonge dans une extension d'une variété abélienne par
un tore.}
\end{enumerate}
\footnote{Un alinéa e), sur le centre de $\pi_1(X)$, est annulé par une série
de boucles.}Le cas c) est un cas particulier de d), puisque le centralisateur
d'un sous-groupe ouvert du $\pi_1$ d'une polycourbe hyperbolique est trivial
(corollaire géométrique de la page 14)\footnote{En marge : « il faut quand
même traiter à part le cas des $\pi_1$ d'une courbe complète de genre
$g \ge 2$, et vérifier que son $\pi_1$ est à centre trivial » — ce qui est
vrai. Pour que c) relève de d), il faut aussi le plongement de $Y$ ; voir la
note de la page 25.}. Dans le cas d), on peut remplacer $X$ par un ouvert,
la normalité assurant que $\pi_1(X') \to \pi_1(X)$ reste surjectif.

\emph{Démonstration.} La situation $(X, Y, f, g)$ provient, par extension des
scalaires, d'une situation analogue sur un corps $K$ de type fini sur
$\mathbb{Q}$, avec $k \supset \overline{K}$. Les groupes $\pi_1$ de $X_K$,
$Y_K$ sont des extensions de $G_K$ par $\Delta_X$, $\Delta_Y$ ; et comme
l'image de $\Delta_X$ a un centralisateur trivial, ces extensions et les
homomorphismes $\pi_1(f_K)$, $\pi_1(g_K)$ sont reconstruits à partir de leurs
parties géométriques et des actions extérieures de $G_K$ (c'est l'argument de
la scholie). Ils sont donc égaux, et tout revient au

\textbf{Corollaire 3.} \emph{Soient $X$, $Y$ de type fini sur un corps $K$
de type fini sur $\mathbb{Q}$, géométriquement connexes, $X$ réduit, $Y$
plongé dans une extension d'une variété abélienne par un tore ; soient $f,
g : X \rightrightarrows Y$ tels que $\pi_1(f) = \pi_1(g)$ comme homomorphismes
d'extensions de $G_K$, à conjugaison par $\Delta_Y$ près. Alors $f = g$.}

Il suffit même que les homomorphismes induits sur les extensions
« abélianisées » — de $G_K$ par $\Delta_X^{\mathrm{ab}}$, $\Delta_Y^{\mathrm{ab}}$
— soient égaux. Une fois les hypothèses anabéliennes ramenées, dit-il, à des
hypothèses « arithmético-géométriques », on peut laisser tomber l'anabélien et
jouer sur l'abélien, mais dans un contexte arithmétique : au lieu de
l'égalité des $H_1(\overline{f})$, $H_1(\overline{g})$ du corollaire 1 a), on a
celle de deux homomorphismes d'\emph{extensions} de $G_K$ par les
$H_1(\overline{X}, \mathbb{Z}_\ell)$, pour tout $\ell$.

\emph{Démonstration.} Il suffit de voir que $f(x) = g(x)$ pour $x \in X(K)$,
car on peut faire une extension finie de $K$ et les points fermés sont denses
dans $X$ réduit. Un point $x$ définit une classe de conjugaison de sections de
$\pi_1(X) \to G_K$, et a fortiori de l'extension abélianisée ; l'hypothèse
dit que $f(x)$ et $g(x)$ définissent la même classe de sections de
$\pi_1(Y)$. Tout revient donc au

\textbf{Théorème 3.} \emph{Soit $X$ un schéma de type fini, géométriquement
connexe, sur un corps $K$ de type fini sur $\mathbb{Q}$, qui se plonge dans
une extension $G$ d'une variété abélienne par un tore. Alors les applications
canoniques}
\[
(\ast)\qquad X(K) \longrightarrow \bigl\{\text{sections de } \pi_1(X) \to
G_K\bigr\}/\Delta_X ,
\]
\[
(\ast\ast)\qquad X(K) \longrightarrow \bigl\{\text{sections de }
\pi_1(X)^{(\mathrm{ab})} \to G_K\bigr\}/\text{conj.}
\]
\emph{sont injectives}, où $\pi_1(X)^{(\mathrm{ab})}$ est l'extension de $G_K$
par $\Delta_X^{\mathrm{ab}} = H_1(\overline{X}, \widehat{\mathbb{Z}})$
déduite de $\pi_1(X)$\footnote{La page donne pour exemple « une variété
élémentaire d'Artin à fibres génériques distinctes de la droite projective et
de la droite affine » ; pour une courbe c'est juste, en dimension $\ge 2$ voir
la note de la page 25. Une note en marge observe qu'il suffirait que $K$ soit
de type fini sur son corps premier, la forme étant « pas très utilisable » en
caractéristique $p$.}.

\emph{Démonstration.} Il suffit de traiter $(\ast\ast)$, et l'on se ramène à
$X = G$. L'application est alors un homomorphisme
\[
G(K) = H^0(K, G) \longrightarrow H^1\bigl(K, T(G)\bigr), \qquad T(G) =
H_1(\overline{G}, \widehat{\mathbb{Z}}) = \varprojlim_n {}_nG .
\]
La suite de Kummer $0 \to {}_nG \to G \xrightarrow{n} G \to 0$ donne
\[
0 \longrightarrow G(K)/n \longrightarrow H^1(K, {}_nG) \longrightarrow
{}_nH^1(K, G) \longrightarrow 0,
\]
et à la limite, l'application $G(K) \to \varprojlim_n H^1(K, {}_nG)$ —
composée de la précédente avec $H^1(K, T(G)) \to \varprojlim H^1(K, {}_nG)$ —
a pour noyau $\bigcap_n nG(K)$, les éléments infiniment divisibles de
$G(K)$\footnote{C'est, sur un schéma en groupes quelconque, la suite de Kummer
que le dossier 141 écrit pour $\mathbb{G}_m$ (pages 3 et 5).}. Or il n'y en a
pas d'autre que $0$\footnote{La page invoque le théorème de Mordell–Weil :
« $G(K)$ est un $\mathbb{Z}$-module de type fini ». C'est vrai pour une
variété abélienne (Lang–Néron), faux dès qu'il y a un tore : $\mathbb{G}_m(K)
= K^{*}$ n'est pas de type fini. Ce qui est vrai, et suffit, est que
$\bigcap_n nG(K) = 0$ : pour la partie abélienne par Lang–Néron ; pour un
tore déployé parce que $K^{*}$ est produit d'un groupe fini (les racines de
l'unité de $K$) et d'un groupe abélien libre, $K$ étant de type fini ; et un
élément de $G(K)$ infiniment divisible a une image infiniment divisible dans la
partie abélienne, nulle, puis (quitte à multiplier par l'exposant de la
torsion de celle-ci) appartient à $\bigcap_n nT(K) = 0$.}, d'où l'injectivité.

\emph{Remarque} (page 33). Il ne suffisait pas de prendre l'extension par
$H_1(\overline{X}, \mathbb{Z}_\ell)$ pour un seul $\ell$ : la translation par
un point d'ordre $p \ne \ell$ donne un contre-exemple. L'énoncé dit ceci :
pour $x \in G(K)$, le torseur sous $T(G)$ image inverse de $x$ dans le
« revêtement universel abélien » $\widetilde{G} = \varprojlim (G
\xrightarrow{n} G)$ est trivial — c'est-à-dire que $x$ est divisible par tout
$n$ — seulement si $x = 0$.

\pagerange{33}{35}
\subsection*{En caractéristique $p$ (pages 33 à 35)}

L'énoncé reste vrai en caractéristique $p > 0$ si l'on remplace les ${}_pG$,
non étales, par des quotients convenables, dans une catégorie de
« quasi-variétés » où les morphismes finis radiciels surjectifs sont inversés.
Mais cette forme « me fait une belle jambe » : il faudrait la version
première à $p$ ; or déjà pour une courbe elliptique, l'injectivité de
$(\ast\ast)$ est fausse en version première à $p$ (translation par un point
d'ordre $p$). Il est possible, ajoute-t-il, que $(\ast)$ reste injective sans
hypothèse de caractéristique — « à examiner de plus près »\footnote{Ces deux
pages sont, dans la transcription, les plus lacunaires du §1 : une dizaine de
segments illisibles. On ne rend que ce qui se lit avec sûreté.}.

Revenant au corollaire 1 b) en caractéristique $p$, avec les groupes
fondamentaux premiers à $p$, la démonstration arithmétique montre que le $u$
de la page 25 est un point de $p$-torsion de $G(K)$ ; si $G$ est un tore (fibres
de genre $0$), $u = 0$. Dans le cas général elle donne seulement : si $P =
f(x)$ et $Q = g(x)$, $Q - P$ est de $p$-torsion dans la jacobienne généralisée
de $Y$, et de même au-dessus de tout revêtement étale galoisien fini ; « il se
pourrait qu'on puisse en déduire $P = Q$ ». En marge : le théorème 2 reste
certainement vrai en caractéristique $p$ pour les groupes premiers à $p$, la
réduction géométrique ramenant au cas d'un produit de courbes de genre
zéro.

\pagerange{36}{39}
\subsection*{Complément : une démonstration géométrique (pages 36 à 39)}

Grothendieck revient au corollaire 1 b) par voie géométrique, en prenant pour
$G$ la jacobienne généralisée de $Y$, et énonce :

\textbf{Lemme.} \emph{Soient $Y$ une polycourbe hyperbolique sur $k$
algébriquement clos, $Y \hookrightarrow J^1_Y$ son plongement dans sa
jacobienne généralisée, $u \in J^0_Y(k)$, et $V \subset Y$ un ouvert non vide
tel que $V + u \subset Y$. Alors $u = 0$.}

Tel quel, sans hypothèse sur les $\pi_1$, le lemme est faux pour les courbes de
types $(1,1)$ et $(1,2)$\footnote{Exemples, sur une courbe elliptique $E$
d'origine $o$. Type $(1,1)$ : $Y = E \setminus \{o\}$, dont la jacobienne
généralisée est $E$ ; pour $u \ne 0$, $V = E \setminus \{o, -u\}$ vérifie $V +
u = E \setminus \{u, o\} \subset Y$. Type $(1,2)$ : $Y = E \setminus \{o,
t\}$ avec $t$ d'ordre $2$ ; la translation par $t$ est un automorphisme de $Y$.
Dans les deux cas l'hypothèse de ($\ast$), page 26, est violée : dans le
premier, le lacet autour de $-u$ est tué par $j$ et envoyé par $j'$ sur un
lacet autour de la pointe $o$, non trivial ; dans le second, $\tau_t$ échange
les deux classes de lacets de pointes. Le corollaire 1 b) n'est donc pas mis en
cause, seulement le lemme du Complément, qui oublie l'hypothèse.} ; il est
vrai pour les autres types de courbes hyperboliques, et pour tous si l'on
garde l'hypothèse de ($\ast$). C'est ce que l'argument de la page donne, une
fois remis en ordre. Par dévissage on se ramène au cas d'une courbe, $Y = X
\setminus I$, $X$ propre de genre $g$, $I$ fini, $V = X \setminus J$ avec $J
\supset I$\footnote{Il note $\hat Y$ la complétion de $Y$ ; nous la notons
$X$.}.

\emph{Genre $g \ge 2$.} Si $Y$ est complète, $V + u \subset Y$ entraîne $Y + u
= Y$, et $\tau_u$ induit un automorphisme de $Y$ qui agit trivialement sur la
cohomologie (il est homotope à l'identité dans $J^1_Y$, dont la cohomologie
s'envoie isomorphiquement sur celle de $Y$ en degré $1$). Par la formule des
points fixes de Lefschetz il a $2 - 2g \ne 0$ points fixes comptés avec
multiplicité, alors qu'une translation non nulle n'en a aucun : $u = 0$. En
général on obtient ainsi que l'image $\bar u$ de $u$ dans $J^0_X$ est nulle ;
et alors $x + u = y$ avec $x, y \in Y$ force $x$ et $y$ à avoir même image
dans $J^1_X$, donc $x = y$ dès que $g \ge 1$, et $u = 0$\footnote{Un
passage barré (pages 36-37) faisait le même calcul sur la courbe « pincée »
$X/I$, en décrivant $J^0_Y$ comme extension de $J^0_X$ par
$\mathbb{G}_m^I/\mathbb{G}_m$. Grothendieck note entre crochets que « ces
arguments ne marchent que pour $g \ne 1$ » et passe à la forme
cohomologique.}.

\emph{Forme cohomologique.} Soient $i, j : V \rightrightarrows Y$ l'inclusion
et $y \mapsto y + u$. On a $H_1(i) = H_1(j)$ : composées avec l'isomorphisme
$H_1(Y) \simeq H_1(J^1_Y)$, elles diffèrent par $H_1(\tau_u) = \mathrm{id}$.
Soit $f : X \xrightarrow{\sim} X$ le prolongement de $j$ aux complétées ; on a
$f(V) \subset Y$, donc $J \supset f^{-1}(I)$. Dans $H_1(V) $ les sous-groupes de
lacets $L_p$, $p \in J$, engendrent $\widehat{\mathbb{Z}}^J/\widehat{\mathbb{Z}}$
(quotient par la diagonale) ; le noyau de $H_1(i)$ est engendré par les $L_p$,
$p \in J \setminus I$, celui de $H_1(j)$ par les $L_p$, $p \in J \setminus
f^{-1}(I)$\footnote{La page écrit $J \setminus f(I)$ ; c'est $f^{-1}(I)$,
puisque $j$ tue le lacet autour de $p$ si et seulement si $f(p) \notin I$.}.
Or dans $\widehat{\mathbb{Z}}^J/\widehat{\mathbb{Z}}$, deux parties propres de
$J$ engendrent le même sous-module si et seulement si elles sont égales ou
toutes deux de cardinal $\operatorname{card} J - 1$. Si $\operatorname{card} I
\ge 2$, les parties $J \setminus I$ et $J \setminus f^{-1}(I)$ sont de cardinal
$\le \operatorname{card} J - 2$, donc égales : $f^{-1}(I) = I$, et $\tau_u$
induit un automorphisme de $Y$\footnote{La page veut écarter le cas
exceptionnel en rapetissant $V$ de sorte que $\operatorname{card}(J \setminus
I) > 1$. Cela ne suffit pas : le cas exceptionnel est $\operatorname{card} I
= 1$, que rapetisser $V$ ne change pas. C'est là que tombent les exemples de
la note précédente.}. Appliquant le même argument à $Y' = X \setminus (I
\setminus \{i\})$ pour chaque $i \in I$ — possible si $\operatorname{card} I
\ge 3$ —, $f$ fixe chaque point de $I$. En genre $0$ ($\operatorname{card} I
\ge 3$ par hyperbolicité) $f$ fixe trois points, donc est l'identité. En genre
$1$ avec $\operatorname{card} I \ge 3$, $f$ est une translation ayant un point
fixe, donc $\bar u = 0$, et l'on conclut comme plus haut. « Ouf ! »

Restent les types $(1,1)$ et $(1,2)$, où l'hypothèse de ($\ast$) est
nécessaire. Elle suffit : si $\pi_1(i)$ et $\pi_1(j)$ coïncident
extérieurement sur un sous-groupe ouvert, un lacet autour de $p \in J$ (ou
l'une de ses puissances) est envoyé par $i$ et par $j$ sur des lacets
conjugués ; comme les sous-groupes d'inertie de deux pointes distinctes de
$Y$ n'ont pas de conjugués commensurables, et que $L_q$ est non trivial dans
$\pi_1(Y)$, $f$ fixe chaque pointe ; en genre $1$ c'est une translation à
point fixe, $\bar u = 0$, et $u = 0$\footnote{Ce paragraphe est de cette
édition. Pour le cas $g = 1$, une note en marge renvoyait déjà à des
« arguments cohomologiques plus bas ».}.

\pagerange{40}{51}
\section*{§2. La question de pleine fidélité (pages 40 à 51)}

\pagerange{40}{44}
\subsection*{Homomorphismes admissibles, et la question sur les corps (pages 40 à 44)}

Soient $K$, $K'$ de type fini sur $\mathbb{Q}$. Tout homomorphisme
$\Pi_{K'} \to \Pi_K$ au-dessus de $\Pi_{\mathbb{Q}}$ provient-il d'un
plongement $K \to K'$ ? On se ramène au cas où $K$ et $K'$ ont la même
fermeture algébrique $k$ de $\mathbb{Q}$, plongée dans
$\overline{\mathbb{Q}}$ ; la question est alors de savoir si tout
homomorphisme $\Delta_{K'} \to \Delta_K$ qui commute à l'action de $G_k$ est
induit par un plongement. Il faudrait, pour en construire un, savoir
reconstruire $K$ à partir de l'extension $G_K$, et l'on « pressent » que le
théorème 3, appliqué à $\mathbb{P}^1_K$ convenablement troué, pourrait en
donner la clé. L'idée : les opérations de $G_k$ sont « si draconiennes »
qu'un homomorphisme qui les respecte ne peut être que géométrique — et cela
exige que le corps de base soit $\mathbb{Q}$, ou de type fini sur
$\mathbb{Q}$, et non quelconque.

Il faut se borner aux homomorphismes d'image ouverte : sinon $K' = k$,
$\Delta_{K'} = 1$, donnerait un plongement $K \to k$. Pour travailler avec les
$\Delta$ plutôt qu'avec les $G$, il faut au moins que le centralisateur dans
$\Delta_K$ de l'image de tout sous-groupe ouvert soit trivial ; un tel
homomorphisme est dit \emph{admissible}. Pour $K' = k$, ce qui comptera n'est
pas l'homomorphisme trivial $1 \to \Delta_K$, mais les \emph{sections} de $G_K
\to G_k$.

\textbf{Question-conjecture.} \emph{Pour un morphisme $B_{K'} \to B_K$ de
topos au-dessus de $B_{\mathbb{Q}}$, les conditions suivantes sont-elles
équivalentes : a) il provient d'un plongement $K \subset K'$ ; b)
l'homomorphisme extérieur $G_{K'} \to G_K$ a une image ouverte ; [c) il est
admissible] ?}

On a a) $\Rightarrow$ b) $\Rightarrow$ c), et b) équivaut à ce que $\Delta_{K'}
\to \Delta_K$ ait une image ouverte\footnote{La page écrit ici $\pi_K \to
\pi_{K'}$, à l'inverse des autres flèches de la page ; nous suivons le sens
$K' \to K$. La condition c) est mise entre crochets, avec en marge « c) n'est
pas juste, cf. plus bas ».}. Une réponse affirmative avec c) impliquerait,
pour $\operatorname{degtr} K' < \operatorname{degtr} K$, qu'aucun homomorphisme
admissible $G_{K'} \to G_K$ n'existe ; en particulier qu'une section de $G_K$
au-dessus d'un ouvert $\Gamma'$ de $G_k$ a toujours un centralisateur non
trivial, donc (le centralisateur de $\Gamma'$ dans $G_k$ étant trivial)
que $\Delta_K^{\Gamma'} \ne 1$. Or cela est « sans doute faux » (le §3 dira
pourquoi : pour les sections venant de points, on attend des invariants
triviaux). D'où la correction :

\begin{enumerate}
\item[c$'$)] \emph{l'homomorphisme $G_{K'}^{\circ} \to G_K$ induit est
admissible.}
\end{enumerate}

Pour voir que c$'$) est nécessaire, il faudrait savoir que pour tout ouvert
$E' \subset G_K$, le centralisateur de $E'^{\circ}$ est trivial ; cela
résulterait de la démonstration du théorème 1 et du fait que, pour $\Gamma'$
ouvert dans $G_{\mathbb{Q}}$, le centralisateur de $\Gamma'^{\circ}$ dans
$G_{\mathbb{Q}}$ est trivial\footnote{En marge : « à vérifier ! ». Nous ne le
vérifions pas ; c'est la seule chose dont dépend la nécessité de c$'$).}.

\textbf{Conséquence (conjecturale).} \emph{Si $K$ n'est pas algébrique sur
$\mathbb{Q}$, pour toute section de $G_K$ au-dessus d'un ouvert $\Gamma'$ de
$G_{\mathbb{Q}}$, on a $\Delta_K^{\Gamma'^{\circ}} \ne 1$.}

Ce n'est pas contradictoire avec ce qu'on attendra des modèles (page 48) :
une section du groupe d'un \emph{corps} normalise, d'après les conjectures
du §3, un sous-groupe d'inertie $\simeq \widehat{\mathbb{Z}}(1)$, sur lequel
$G^{\circ}$ agit trivialement — c'est ce que les pages 55 et 60 rendront
explicite\footnote{Ce rapprochement est le nôtre ; Grothendieck annonce
seulement, page 48, qu'il étudiera « les relations entre cette conséquence
conjecturale et la précédente (d'apparence opposée !) ».}.

\pagerange{45}{49}
\subsection*{Modèles, et la conjecture fondamentale (pages 45 à 49)}

Les $G_K$ sont, à certains égards, trop gros ; il faut les regarder comme
limites des $\pi_1(U)$ de leurs modèles affines, et particulièrement de
modèles élémentaires — ou plus généralement de modèles qui sont des
$K(\pi, 1)$ au sens profini. Peut-être faut-il exiger, dans la conjecture,
qu'un homomorphisme $G_{K'} \to G_K$ respecte les « filtrations
modéliques » de ces groupes. On étudie donc la situation analogue pour les
modèles, sur $\mathbb{Q}$ (la situation « mixte », sur un corps de type fini,
lui paraît « un peu bâtarde »). Un morphisme $V \to U$ donne un morphisme de
topos galoisiens $B_V \to B_U$ ; si $U$ est élémentaire hyperbolique, il est
connu dès qu'on connaît son effet sur les $H_1(\overline{V}, \mathbb{Z}_\ell)
\to H_1(\overline{U}, \mathbb{Z}_\ell)$, ce qui est bien moins que la classe
du morphisme de topos\footnote{Le sens des flèches de la page ($B_U \to B_V$,
$E_U \to E_V$) est inverse de celui de $V \to U$ ; nous suivons $V \to U$.}.
Mais quels homomorphismes viennent de morphismes ? « Avec un peu de culot, on
dirait » :

\textbf{Conjecture fondamentale.} \emph{Soient $U$, $V$ des schémas
réguliers de type fini sur $\mathbb{Q}$, $U$ une polycourbe hyperbolique sur
une extension finie de $\mathbb{Q}$, et $f : \pi_1(V) \to \pi_1(U)$ un
homomorphisme extérieur compatible aux projections sur $G_{\mathbb{Q}}$.
Équivalence de : a) $f$ provient (à isomorphisme près) d'un morphisme $V \to
U$, alors unique ; b) $f|_{\pi_1(V)^{\circ}}$ est admissible.}

La nécessité de b) se ramène au cas où $V$ est un point, et c'est la

\textbf{Conséquence conjecturale.} \emph{Si $\Gamma' \subset G_{\mathbb{Q}}$
est ouvert, de corps fixe $k'$, et $x$ un $k'$-point de $U$, d'où une
section $\Gamma' \to \pi_1(U)$ et une action de $\Gamma'$ sur $\Delta_U$,
alors $\Delta_U^{\Gamma'^{\circ}} = 1$.}

La conjecture sur les modèles implique celle sur les corps, si l'on s'y borne
aux homomorphismes compatibles aux filtrations modéliques\footnote{Une note
en marge conseille de « se borner à l'équivalence de a) et b) » ; la condition
c) qu'elle vise n'est pas écrite page 47, c'est vraisemblablement celle des
pages 43-44.}. Plus généralement, pour des schémas $V$ qui sont limites projectives de
schémas réguliers de type fini, à transitions affines (des immersions
ouvertes), par exemple les spectres d'anneaux locaux réguliers
$\mathcal{O}_{X,x}$, les morphismes \emph{dominants} $V \to U$ devraient
correspondre exactement aux homomorphismes extérieurs compatibles à
$G_{\mathbb{Q}}$ et d'image ouverte\footnote{Quand $U$ est une courbe, cette
forme « dominante » est aujourd'hui un théorème : Mochizuki (1999) a démontré
que, pour $X$ lisse et $Y$ courbe hyperbolique sur un corps sous-$p$-adique,
les morphismes dominants $X \to Y$ correspondent aux homomorphismes ouverts
$\pi_1(X) \to \pi_1(Y)$ au-dessus de $G_K$, à conjugaison près. La forme
« isomorphismes » est due, pour les courbes, à Nakamura, Tamagawa et
Mochizuki, et pour les corps de fonctions à Pop.
La forme générale, avec b), contient la conjecture des sections (cas où $V$ est
un point), qui reste ouverte.}.

\pagerange{49}{51}
\subsection*{Un programme de reconstruction (pages 49 à 51)}

La conjecture admise, la question suivante est de déterminer quels topos
galoisiens sur $B_{\mathbb{Q}}$ viennent de modèles élémentaires
hyperboliques — c'est-à-dire, les groupes $\Delta_U$ étant connus, quelles
actions extérieures d'ouverts de $G_{\mathbb{Q}}$ sur eux sont
géométriques ; et la même question pour d'autres $K(\pi, 1)$ réalisés par des
variétés qui ne sont pas élémentaires, au premier chef les variétés
modulaires des courbes algébriques. À partir de là on reconstruirait, par
recollement et en termes profinis, la catégorie des schémas lisses sur un corps
de type fini, puis celle des schémas localement de type fini, au moins à
homéomorphisme universel près.

D'où un énoncé : pour $S$ localement noethérien, le foncteur $X \mapsto
X_{\mathrm{ét}}$, des schémas réduits localement de présentation finie sur $S$
vers les topos au-dessus de $S_{\mathrm{ét}}$, serait fidèle (deux morphismes
$f, g$ tels que $f_{\mathrm{ét}}$ et $g_{\mathrm{ét}}$ soient isomorphes sont
égaux), et peut-être pleinement fidèle après inversion des homéomorphismes
universels\footnote{En marge : « $S$ de car. $0$ ? ». Sur un corps de type
fini de caractéristique $0$, la forme « isomorphismes » pour les schémas normaux
a été démontrée par Voevodsky (1990), et une forme pleinement fidèle, à
homéomorphisme universel près, beaucoup plus récemment (Carlson, Haine et
Wolf) — sauf erreur de notre part sur la portée exacte de ces énoncés.}. Pour
une courbe propre sur un corps de nombres, on retrouverait que tout
automorphisme extérieur de $G_K$ ($K$ son corps de fonctions) qui respecte la
structure et commute à $G_k$ provient d'un automorphisme de la courbe.

\pagerange{52}{103}
\section*{§3. Étude des sections de $\pi_1(U)$ sur $G_K$ (pages 52 à 103)}

\pagerange{52}{55}
\subsection*{Points et pointes (pages 52 à 55)}

Soit $U$ un schéma lisse de type fini géométriquement connexe sur $K$, et la
suite exacte $1 \to \Delta_U \to \pi_1(U) \to G_K \to 1$, définie par le choix
d'un revêtement universel $\widetilde{U}$. On étudie ses sections à
$\Delta_U$-conjugaison près, et en même temps les sections au-dessus d'ouverts
$G_{K'}$ de $G_K$ ($K'$ fini sur $K$). Si $U$ se plonge dans un schéma en
groupes commutatif convenable, $U(K) \to \{\text{sections}\}/\Delta_U$ est
injective (théorème 3). Il y a d'autres sections, « géométriques » elles
aussi.

Soit $U$ une courbe de type $(g, \nu) \ne (0,0), (0,1)$, et $i \in X
\setminus U$ une pointe rationnelle sur $K$. Le sous-groupe de lacets $L_i
\simeq \widehat{\mathbb{Z}}(1)$ a pour normalisateur dans $\pi_1(U)$ le
\emph{groupe de décomposition} de la pointe, extension
\[
1 \longrightarrow L_i \longrightarrow N_{\pi_1(U)}(L_i) \longrightarrow G_K
\longrightarrow 1 ,
\]
et ses scindages donnent des sections de $\pi_1(U)$\footnote{La page écrit
« son centralisateur $Z(L_i)$ ». C'est le normalisateur : $G_K$ agit sur $L_i
\simeq \widehat{\mathbb{Z}}(1)$ par le caractère cyclotomique, donc ne le
centralise pas, et la page 79 dira d'ailleurs « le normalisateur de $L_i$ ». La
suite est exacte à gauche parce que $N_{\Delta_U}(L_i) = L_i$, fait classique
pour une courbe hyperbolique, que le dossier 145 utilise aussi (pages 6 à 11)
pour $\mathbb{P}^1 \setminus \{0,1,\infty\}$.}. Il en existe : une
uniformisante en $i$ en définit une — c'est aujourd'hui un \emph{point base
tangentiel} (Deligne, 1989). Les classes de scindages forment un torseur sous
\[
H^1\bigl(G_K, \widehat{\mathbb{Z}}(1)\bigr) \simeq \varprojlim_n H^1(K, \mu_n)
\simeq \varprojlim_n K^{*}/K^{*n} = \widehat{K^{*}},
\]
et s'envoient injectivement dans les classes de sections de
$\pi_1(U)$\footnote{Affirmé sans démonstration. L'injectivité revient à ce que
l'image d'un tel scindage ne normalise aucun autre conjugué de $L_i$, ce que
la relation ($\dagger$) ci-dessous donnerait.}.

\textbf{Proposition.} \emph{Soit $U$ hyperbolique. Les classes de sections
issues des scindages des groupes de décomposition des pointes sont distinctes
de celles des points de $U(K)$ ; et, sauf peut-être pour le type $(0,3)$,
deux pointes distinctes donnent des classes distinctes.}\footnote{La page
suppose seulement $(g, \nu) \ne (0,0), (0,1)$ pour la première assertion, et
l'hyperbolicité pour la seconde. Pour le type $(0,2)$, $U = \mathbb{G}_m$, la
première est fausse : $\Delta_U = L_0 = L_\infty$, le normalisateur de $L_0$
est $\pi_1(U)$ tout entier, et toute section, celles des points comprises, est
« cuspidale ». En marge : « C'est démontré sauf peut-être pour le type $(0,3)$,
qui mérite un examen approfondi ».}

Pour la première assertion on « bouche » le trou $i$ : la section devient,
dans $U' = U \cup \{i\}$, celle du point $i$, distincte de celles des autres
points de $U'$ par le théorème 3 (appliqué à $U'$, de type $(g, \nu - 1)$, qui
se plonge dans sa jacobienne généralisée), a fortiori distincte de celles des
points de $U$. On raisonne de même pour deux pointes en bouchant les deux
trous, sauf pour le type $(0,3)$ qui tombe sur le type $(0,1)$, où l'on n'a pas
l'injectivité. On s'en tirerait si l'on savait que, pour un scindage du groupe
de décomposition de $i$, faisant agir $G_K$ sur $\Delta_U$,
\[
(\dagger)\qquad \Delta_U^{\,G_K^{\circ}} = L_i \qquad (\text{donc }
\Delta_U^{\,G_K} = 1),
\]
« résultat qui me paraît hautement plausible » ; en marge : « à énoncer parmi
les conjectures-clefs »\footnote{La page numérote cette relation (3), numéro
déjà employé à la page 53 pour l'isomorphisme de Kummer.}. La classe de la
section déterminerait alors celle de $L_i$, donc $i$.

\pagerange{55}{57}
\subsection*{La conjecture A, et le passage au corps de fonctions (pages 55 à 57)}

\textbf{Conjecture A.} \emph{Soit $U$ une courbe hyperbolique géométriquement
connexe sur un corps $K$ de type fini sur $\mathbb{Q}$. Toute section de
$\pi_1(U) \to G_K$ est de l'un des deux types précédents : définie par un
point de $U(K)$, ou par un scindage du groupe de décomposition d'une pointe
$i \in X(K) \setminus U(K)$.}

En marge, une précision : $\Delta_U^{G^{\circ}} = 1$ dans le premier cas,
$= L_i$ dans le second. C'est la \emph{conjecture des sections} de
Grothendieck, dans sa forme pour les courbes affines\footnote{Elle figure dans
la lettre à Faltings de juin 1983. À notre connaissance elle reste ouverte
sur les corps de nombres ; sa forme birationnelle, qui suit, a été démontrée
sur les corps $p$-adiques (Koenigsmann, 2005 ; Pop).}.

Passant à la limite sur les modèles — au corps de fonctions $L$ de $U$, et à
$G_L \to G_K$, groupe « à lacets infini » où chaque point de $X$ est une
pointe —, on obtiendrait que toute section de $G_L \to G_K$ provient d'un
scindage du groupe de décomposition d'un point $x \in X(K)$ : c'est la
\emph{conjecture des sections birationnelle}. Les sections se groupent en
paquets, selon la classe de conjugaison du centralisateur de l'image de
$G_K^{\circ}$, et l'on retrouve $X(K)$ — et les $X(K')$, $K'$ fini sur $K$ —
à partir de l'extension $G_L$ de $G_K$, en même temps qu'une façon de
reconstruire les ouverts $U = X \setminus I$. La structure $G_L \to G_K$ est
donc la plus riche, et de loin la plus commode en genre $0$ et $1$, où se
borner aux $U$ de type $(g, \nu)$ hyperboliques « casse le groupe continu
d'automorphismes ». La forme « modélique » revient à la birationnelle si l'on
précise que toute section de $\pi_1(U) \to G_K$ se relève en une section de
$G_L \to G_K$.

\pagerange{57}{62}
\subsection*{Forme géométrique : un point fixe et un seul (pages 57 à 62)}

Soit $\widetilde{U}$ un revêtement universel de $U$, et $X'$ la normalisée de
$X$ dans $\widetilde{U}$ — que Grothendieck s'abstient de noter
$\widetilde{X}$, « vu que justement il n'est pas étale sur $X$ ». Les fibres
de $X'$ au-dessus d'une pointe $i$ sont en bijection avec les sous-groupes de
lacets conjugués à $L_i$ ; $X'$ est donc la réunion de $\widetilde{U}$ et de
l'ensemble des sous-groupes de lacets, qui jouent le rôle de « points à
l'infini » du revêtement universel.

\begin{tikzcd}
U \arrow[r, hook] & X \\
\widetilde{U} \arrow[r, hook] \arrow[u] & X' \arrow[u]
\end{tikzcd}

\footnote{Le diagramme de la page 57 met en regard les lignes $U \subset X$,
$\overline{U} \subset \overline{X}$, $\widetilde{U} \subset X'$ au-dessus de $K
\subset \overline{K}$, par des traits sans flèches. Nous n'en gardons que les
deux lignes extrêmes, les flèches verticales montant de $\widetilde{U}$ et $X'$
vers $U$ et $X$.}Le groupe $\pi_1(U)$ s'interprète comme groupe
d'automorphismes du schéma $\widetilde{U}$ compatibles à son action sur
$\overline{K}$ (la fermeture algébrique de $K$ dans $H^0(\widetilde{U},
\mathcal{O})$). Une section de $\pi_1(U) \to G_K$ est donc une action de
$G_K$ sur $\widetilde{U}$, compatible à son action sur $\overline{K}$ et
continue en un sens convenable.

\textbf{Conjecture B.} \emph{Toute telle action de $G_K$ sur
$\widetilde{U}$ a, dans $X'$, un point fixe et un seul.}

Cela signifie : s'il y a un point fixe $\tilde x$ « à distance finie », son
image dans $U$ est déterminée (théorème 3), il n'y en a pas d'autre au-dessus
($\Delta_U^{G_K} = 1$), et il n'y a pas en même temps de point fixe à l'infini
(établi plus haut) ; s'il y a un point fixe à l'infini, une même action ne
normalise pas deux $L_i$, $L_j$ distincts (établi sauf pour le type $(0,3)$), et
le $L_i$ normalisé est unique — affaiblissement de ($\dagger$). Grothendieck
conjecture davantage : \emph{$G_K^{\circ}$ a déjà, dans $X'$, un point fixe et
un seul} (nécessairement fixe par $G_K$). À distance finie, cela donne
$\Delta_U^{G^{\circ}} = 1$ ; à l'infini, $\Delta_U^{G^{\circ}} \subset
N_{\Delta_U}(L_i) = L_i$, d'où la relation ($\dagger$).

Il se demande alors si ces énoncés ne vaudraient pas en partant seulement
d'un relèvement de $G_K^{\circ}$, ou d'un sous-groupe contenant un ouvert de
$G_K^{\circ}$ ; et s'ils ne vaudraient pas dès que $K$ est de type fini sur un
corps cyclotomique, fini ou non sur $\mathbb{Q}$.

Il appelle \emph{courbe de Poincaré} sur $\overline{K}$ algébriquement clos de
caractéristique $0$ un schéma isomorphe au revêtement universel d'une courbe
hyperbolique sur $\overline{K}$ (deux courbes donnent des courbes de Poincaré
isomorphes si et seulement si elles ont des revêtements finis étales
isomorphes) ; une telle courbe a une complétion canonique, sur laquelle ses
automorphismes agissent.

\textbf{Conjecture B$'$.} \emph{Soient $K$ de type fini sur $\mathbb{Q}$ (ou
sur un corps cyclotomique), $\mathcal{U}$ une courbe de Poincaré sur
$\overline{K}$, de complétion $\mathcal{X}$, et une action de $G_K$ sur
$\mathcal{U}$ compatible à son action sur $\overline{K}$. Alors $G_K^{\circ}$
a dans $\mathcal{X}$ un point fixe et un seul.}

La différence avec B : on ne suppose \emph{pas} que l'action normalise celle
d'un groupe profini $\Delta$ agissant librement avec quotient une courbe
algébrique.

\pagerange{63}{68}
\subsection*{Courbes relatives, et un programme de reconstruction (pages 63 à 68)}

Les conjectures s'appliquent quand $K$ est le corps de fonctions d'un modèle
$S$ (une polycourbe hyperbolique, disons) et que $U_K$ provient d'une courbe
relative $U_S \to S$. On a alors $\pi_1(U_S) \to \pi_1(S)$ de noyau $\Delta$,
et $\pi_1(U_K) \to G_K$ s'en déduit par produit fibré le long de la
surjection $G_K \to \pi_1(S)$ :

\begin{tikzcd}
\pi_1(U_K) \arrow[r] \arrow[d] & \pi_1(U_S) \arrow[d] \\
G_K \arrow[r] & \pi_1(S)
\end{tikzcd}

Les sections de $\pi_1(U_K)$ sur $G_K$ sont donc les relèvements continus $G_K
\to \pi_1(U_S)$ de $G_K \to \pi_1(S)$, et ceux qui sont triviaux sur le noyau de
$G_K \to \pi_1(S)$ sont exactement les sections de $\pi_1(U_S)$ sur $\pi_1(S)$.
Les conjectures prédisent qu'elles sont de deux sortes : 1°) celles des
sections rationnelles de $U_S$ sur $S$ — et l'on vérifiera (pages 69 à 78)
qu'une telle section rationnelle ne donne une section de $\pi_1(U_S)$ que si
elle est \emph{régulière} ; 2°) celles des sections sur $S$ du revêtement fini
étale $S' = X_S \setminus U_S$ des pointes. Restent deux points « un peu
techniques » : à quelle condition une telle section définit-elle des scindages,
et lesquels.

Avant, il veut voir dans quelle mesure la conjecture A permettrait de
reconstruire la catégorie des modèles élémentaires hyperboliques sur
$\mathbb{Q}$ et de leurs corps de fonctions, en termes de groupes extérieurs,
à partir de la donnée fondamentale : \emph{$G_{\mathbb{Q}}$ opérant
extérieurement sur $\widehat{\pi}_{0,3}$}, groupe fondamental de
$\mathbb{P}^1_{\mathbb{Q}} \setminus \{0, 1, \infty\}$, d'où l'extension
$\pi_1(\mathbb{P}^1 \setminus \{0,1,\infty\})$ de $G_{\mathbb{Q}}$\footnote{C'est
l'action étudiée dans le dossier 145, dont les automorphismes « à lacets » de
$\widehat{\pi}_{0,3}$ ont pour multiplicateur précisément le caractère
cyclotomique ; le noyau $G^{\circ}$ de ce dossier-ci est le sous-groupe de
multiplicateur $1$.}. Les classes de sections « admissibles » (centralisateur
de l'image de $G^{\circ}$ trivial) de cette extension forment un ensemble sur
lequel $G_{\mathbb{Q}}$ opère, qui devrait être $(\mathbb{P}^1 \setminus
\{0,1,\infty\})(\overline{\mathbb{Q}})$ à un ensemble près. Pour tout ensemble
fini $I$ de telles sections, stable par $G_{\mathbb{Q}}$, un foncteur de
« forage des trous » doit fournir un groupe extérieur de type $(0, 3 +
\operatorname{card} I)$ sur lequel $G_{\mathbb{Q}}$ opère. Il vaudrait mieux,
dit-il, utiliser « le yoga introduit par Ihara des groupoïdes
rigides »\footnote{La page ne dit pas de quel texte d'Ihara il s'agit, et
nous ne l'identifions pas.}, et construire l'équivalent groupoïdal de
$\mathbb{P}^1_{\mathbb{Q}}$, la « sphère arithmétique », qui correspond à un
groupe à lacets infini — c'est $G_{K_1}$, $K_1 = \mathbb{Q}(T)$. De proche en
proche, par des fibrations dont chaque cran est une courbe de genre $0$ ou un
revêtement étale fini d'une telle, on construirait tous les $G_K$ au-dessus de
$G_{\mathbb{Q}}$ (tout corps de type fini est fini sur une extension
transcendante pure) et tous les modèles élémentaires, assez pour avoir un
système fondamental de voisinages de tout point d'une variété lisse, et
reconstituer celle-ci par recollement.

« Mais je suis certain que pour faire une telle description, il me faudra
développer un langage géométrique qui colle mieux à l'intuition géométrique,
que les sempiternelles extensions de groupes profinis… ou actions
extérieures, et où les points rigides (ou cuspidaux) dans des clôtures
algébriques de corps jouent un rôle prépondérant. »

\pagerange{69}{71}
\subsection*{Sections rationnelles et sections partout définies (pages 69 à 71, et 82-83)}

Soient $S$ lisse et irréductible sur $\mathbb{Q}$, de corps de fonctions $K$,
$U_S \to S$ élémentaire à fibres de $\pi_1$ non nul, et $f$ une section
rationnelle de $U_S$ sur $S$, d'où une section de $\pi_1(U_K)$ sur $G_K$,
c'est-à-dire un relèvement $G_K \to \pi_1(U_S)$. Il veut montrer : \emph{$f$
est partout définie si et seulement si ce relèvement s'annule sur le noyau de
$G_K \to \pi_1(S)$.}

La condition est « de codimension $1$ » : si $Z \subset S$ est fermé de
codimension $\ge 2$ et $S' = S \setminus Z$, on a $\pi_1(S')
\xrightarrow{\sim} \pi_1(S)$ (pureté de Zariski–Nagata), donc les noyaux de
$G_K \to \pi_1(S)$ et $G_K \to \pi_1(S')$ coïncident. Il faut donc s'assurer
qu'une section de $U_S$ définie sur $S'$ se prolonge à $S$. De proche en
proche, on est ramené à une courbe relative $U_S = X_S \setminus T$, $X_S$
lisse et propre de dimension relative $1$, $T$ fini étale sur $S$. Si le
genre est $\ge 1$, la section se prolonge en une section de $X_S$ (théorème de
Weil sur le prolongement des applications rationnelles vers les schémas
abéliens), et l'image inverse de $T$ est un diviseur de $S$ dont la trace sur
$S'$ est nulle, donc nul.

\pagerange{82}{83}
Le feuillet des pages 82 et 83\footnote{Sa page 41. Sur la place de ce
feuillet, voir « L'ordre des feuillets » plus haut.} traite le genre $0$ :
prenant une section de $T$ pour section à l'infini, on peut supposer (en
localisant) $U_S = \mathbb{A}^1_S \setminus T'$, et $f$ est une fonction sur
$S'$, qui se prolonge à $S$ normal. Puis il introduit les groupes de
décomposition et d'inertie aux points de codimension $1$ : si $x$ est le
point générique d'un diviseur irréductible $D$ de $S$\footnote{La page écrit
« les points $x_i$ de $X$ » ; il s'agit de $S$.}, $\mathcal{O}_{S,x}$ est un
anneau de valuation discrète ; on choisit un point $\tilde x$ au-dessus dans la
normalisée de $S$ dans $\overline{K}$, et son stabilisateur $N_{\tilde x}
\subset G_K$ s'envoie sur $\mathrm{Gal}(\overline{k(x)}/k(x))$, avec pour noyau
le groupe d'inertie :
\[
1 \longrightarrow I_{\tilde x} \longrightarrow N_{\tilde x} \longrightarrow
G_{k(x)} \longrightarrow 1, \qquad I_{\tilde x} \simeq
\widehat{\mathbb{Z}}(1),
\]
la dernière identification par Kummer, « à cause de l'hypothèse de
caractéristique $0$ » : l'inertie est alors modérée\footnote{C'est le calcul du
dossier 141, page 1 : sur le corps des fractions d'un anneau de valuation
discrète hensélien à corps résiduel de caractéristique $0$, toute extension
finie est de Kummer.}. Ici $k(x)$, corps des fonctions de $D$, est de type fini
sur $\mathbb{Q}$. Le noyau de $G_K \to \pi_1(S)$ est le sous-groupe distingué
fermé engendré par les $I_{\tilde x}$ (pureté), et la condition envisagée
revient à ce que le relèvement s'annule sur chacun d'eux\footnote{La fin de la
page 83, deux lignes rapides, ne se lit pas avec assurance.}.

\pagerange{72}{78}
\subsection*{Le critère local (pages 72 à 78)}

On passe à l'anneau local $\mathcal{O}_x$ en un tel $x$ : le relèvement $G_K
\to \pi_1(U_S)$ se factorise par un relèvement $\varphi_x : G_K \to
\pi_1(U_{\mathcal{O}_x})$ de $G_K \to \pi_1(\operatorname{Spec}
\mathcal{O}_x)$, et il s'agit de prouver que $\varphi_x$ s'annule sur
$I_{\tilde x}$ si et seulement si $f$ est définie en $x$\footnote{Une note
marginale oblique, en grande partie douteuse, observe que le quotient de $G_K$
par le sous-groupe distingué engendré par $I_x$ est
$\pi_1(\operatorname{Spec} \mathcal{O}_x)$.}. C'est un énoncé purement
géométrique, qu'il formule en général :

\textbf{Théorème} (page 73)\textbf{.} \emph{Soient $S$ un trait (spectre
d'un anneau de valuation discrète) de corps des fractions $K$ et de point
fermé $s$, et $U \to S$ une courbe relative élémentaire (plus généralement un
schéma élémentaire dont les fibrations n'ont pas de fibres de type $(0,0)$ ou
$(0,1)$). Soit $f_K \in U(K)$, d'où une section $\psi$ de $\pi_1(U_K)$ sur $G_K$.
Équivalence de : a) $f_K$ se prolonge en une section de $U$ sur $S$ ; b)
$\psi$ provient d'une section de $\pi_1(U)$ sur $\pi_1(S)$ ; b$'$) le composé
$G_K \to \pi_1(U_K) \to \pi_1(U)$ s'annule sur l'inertie $I_s$.}\footnote{La
page numérote ce théorème « 1) », lecture incertaine, et écrit $T$ puis $S$
pour le même trait.}

b) $\Leftrightarrow$ b$'$) et a) $\Rightarrow$ b) sont clairs ; reste b)
$\Rightarrow$ a). Par descente on suppose $S$ strictement local (alors $G_K =
I_s$ et $\pi_1(S) = 1$) ; de proche en proche $U = X \setminus T$ est une
courbe, $X$ lisse et propre, et $f$ se prolonge en une section de $X$. Il faut
voir que $f(S) \subset U$. D'où le

\textbf{Lemme.} \emph{Soient $S$ un trait strictement local, $X \to S$ propre et
lisse de dimension relative $1$, à fibres connexes, $T = \bigsqcup_{i \in I}
g_i(S)$ réunion de sections disjointes, $U = X \setminus T$, de type $(g,
\nu) \ne (0,1)$. Soient $i_0 \in I$ et $f$ une section de $X$, distincte des
$g_i$, avec $f(s) = g_{i_0}(s)$. Si $Y = S \setminus \{s\}$, l'homomorphisme
$\pi_1(Y) \to \pi_1(U)$ induit par $f$ n'est pas trivial ; et même, pour $\nu
\ge 2$ et $\ell$ premier à la caractéristique résiduelle,}
\[
\mathbb{Z}_\ell(1) \simeq H_1(Y, \mathbb{Z}_\ell) \longrightarrow H_1(U,
\mathbb{Z}_\ell) \simeq H_1(U_s, \mathbb{Z}_\ell)
\]
\emph{n'est pas nul.}

Comme $f$ définit une section régulière du torseur jacobien $J^1_{X/S}$,
l'image tombe dans la partie torique de $H_1(U, \mathbb{Z}_\ell)$, isomorphe à
$\mathbb{Z}_\ell(1)^I/\mathbb{Z}_\ell(1)$ (nulle si $\nu = 1$, auquel cas le
critère homologique ne suffit pas). Il faut calculer
\[
\mathbb{Z}_\ell(1) \longrightarrow \mathbb{Z}_\ell(1)^I/\mathbb{Z}_\ell(1),
\]
et il « devine » le résultat\footnote{Un calcul intrinsèque, par les isogénies
de $J^1_{U/S}$ et $\mathrm{Ker}(J^1_{U/S} \to J^1_{X/S}) \simeq
\mathbb{G}_m^I/\mathbb{G}_m$, est commencé puis barré de trois grands
traits.} : soient $A = \mathcal{O}_{X, g_{i_0}(s)}$, $V$ l'anneau de $S$,
$J_{i_0}$ et $J_f$ les idéaux (inversibles) des sections $g_{i_0}$ et $f$, et
\[
n = \operatorname{long}\, V/g_{i_0}^{*}(J_f) = \operatorname{long}\,
A/(J_{i_0} + J_f),
\]
la multiplicité d'intersection de $f(S)$ et $g_{i_0}(S)$, finie puisque $f \ne
g_{i_0}$. L'homomorphisme serait $n$ fois l'injection canonique d'indice
$i_0$. « Il faudrait que je demande à Carlos de me le confirmer. »\footnote{Le
prénom est tel qu'écrit ; nous n'identifions pas la personne. La présomption
est juste : si $t$ est une équation locale de $g_{i_0}(S)$ et $\varpi$ une
uniformisante de $V$, $f^{*}t = (\text{unité}) \cdot \varpi^{n}$, et le lacet
autour de $s$ s'envoie sur $n$ fois le lacet autour de la pointe $i_0$ dans
le groupe fondamental modéré du voisinage épointé ; c'est vrai au niveau des
$\pi_1$ (premiers à la caractéristique résiduelle), pas seulement des $H_1$.}

Pour $\nu = 1$ il faut vérifier au niveau des $\pi_1$ (premiers à $p$) que
$\pi_1(Y) \to \pi_1(U_s)$ est $k_{i_0} \circ (n \cdot \mathrm{id})$, où $k_{i_0}
: \widehat{\mathbb{Z}}(1) \to \pi_1(U_s)$ est l'homomorphisme de lacets local
de la pointe $i_0$ — « je vais admettre ce point, qui ne peut guère être
difficile » ; c'est le même calcul local que pour $\nu \ge 2$, et l'énoncé, comme
le dit une note en marge, est indépendant de $\nu$. Comme $L_{i_0}$ est non
trivial et sans torsion dans $\Delta_U$, b) $\Rightarrow$ a) en
résulte.

\pagerange{78}{84}
\subsection*{Sections de seconde espèce, et la classe conormale (pages 78 à 81, puis 84)}

Soient toujours $U = X \setminus T$ sur $S$, $X$ propre et lisse, $T$ fini
étale, et $g_i$ une section de $T$. On suppose la suite $1 \to \Delta \to
\pi_1(U) \to \pi_1(S) \to 1$ exacte — ce qui demande que $\pi_2(S) \to
\pi_1(\text{fibre})$ soit nul, vrai par exemple si $S$ est un modèle élémentaire
en caractéristique $0$\footnote{Une longue note en marge, dont un premier bloc
est barré, observe que l'action extérieure de $\pi_1(S)$ sur le $\pi_1$ de la
fibre géométrique suffit à construire une extension, et conclut « Il faudrait
regarder ceci de plus près ! ». La place de son astérisque dans le texte n'a
pas été retrouvée.}. Le normalisateur de $L_i$ dans $\pi_1(U)$ s'envoie
\emph{sur} $\pi_1(S)$ et donne une extension
\[
1 \longrightarrow \widehat{\mathbb{Z}}(1) \longrightarrow N(L_i)
\longrightarrow \pi_1(S) \longrightarrow 1,
\]
dont la classe est un élément $c_i \in H^2(\pi_1(S), \widehat{\mathbb{Z}}(1))
\simeq H^2(S, \widehat{\mathbb{Z}}(1))$ si $S$ est un $K(\pi, 1)$. La suite de
Kummer
\[
0 \to \operatorname{Pic}(S)/n \to H^2(S, \mu_n) \to {}_nH^2(S, \mathbb{G}_m) \to 0
\]
donne à la limite $0 \to \widehat{\operatorname{Pic}(S)} \to H^2(S,
\widehat{\mathbb{Z}}(1)) \to \varprojlim_n {}_n\mathrm{Br}(S)$. Si $S$ est
élémentaire sur un corps de type fini sur $\mathbb{Q}$, $\operatorname{Pic}(S)$
est de type fini (Mordell–Weil–Néron), et $\operatorname{Pic}(S) \to H^2(S,
\widehat{\mathbb{Z}}(1))$ est injectif.

Il affirme que $c_i$ est l'image de la classe $\xi_i \in \operatorname{Pic}(S)$
du fibré conormal — « ou normal, le diable sait » — de $g_i(S)$ dans $X$, et
esquisse la vérification : sur le complété formel de $X$ le long de $g_i(S)$ —
un voisinage tubulaire —, l'extension est la suite d'homotopie de ce
« fibré en disques épointés » au-dessus de $S$\footnote{Les deux signes
diffèrent par $\xi_i \mapsto -\xi_i$, et le choix dépend des conventions de la
suite de Kummer ; nous ne le tranchons pas, la page non plus. La vérification
n'est pas faite dans le dossier.}. Dans le cas « arithmétique »,
l'extension est donc scindée si et seulement si $\xi_i = 0$, c'est-à-dire s'il
existe globalement sur $S$ une uniformisante de $X$ le long de $g_i(S)$ —
condition toujours satisfaite localement, et suffisante en général.

Quand $\xi_i = 0$, parmi les scindages il y a ceux qui proviennent d'une base
de $J_i/J_i^2$ ; leur indétermination est dans $\mathbb{G}_m(S)$, celle de tous
les scindages dans
\[
H^1(S, \widehat{\mathbb{Z}}(1)), \qquad 0 \to \widehat{\mathbb{G}_m(S)} \to H^1(S,
\widehat{\mathbb{Z}}(1)) \to \varprojlim_n {}_n\operatorname{Pic}(S),
\]
et dans le cas arithmétique le dernier terme est nul (module de Tate d'un
groupe de type fini), de sorte que $H^1(S, \widehat{\mathbb{Z}}(1)) \simeq
\widehat{\mathbb{G}_m(S)}$\footnote{Le premier terme de la suite de Kummer de la
page 81 est surchargé ; la lecture $\mathbb{G}_m(S)/\mathbb{G}_m(S)^n$ est celle
qu'exige la limite qui suit.}. Comme $\mathbb{G}_m(S)$ n'a pas d'élément
infiniment divisible autre que $1$ (page 84), $\mathbb{G}_m(S) \to
\widehat{\mathbb{G}_m(S)}$ est injectif : les scindages « géométriques »
forment un torseur sous $\mathbb{G}_m(S)$, sous-torseur du torseur sous
$\widehat{\mathbb{G}_m(S)}$ de tous les scindages. « Pour que la description
profinie de la géométrie algébrique absolue sur $\mathbb{Q}$ soit complète, il
y faudrait également caractériser » ce sous-ensemble remarquable — ce qui
n'est pas fait.

\pagerange{84}{98}
\subsection*{Spécialisation : d'une section de première espèce à une de seconde (pages 84 à 98)}

Comment une section de première espèce peut-elle se « spécialiser » en une de
seconde espèce ? Soit $f$ une section rationnelle de $U$ sur $S$, qui en un
point $s$ de codimension $1$ coïncide, comme section de $X$, avec une $g_i$ :
$f(s) = g_i(s)$. Alors $G_K \to \pi_1(U_K) \to \pi_1(U)$ est non trivial sur
l'inertie $I_{s'}$, qu'il envoie dans $\Delta$, et même dans l'un des $L_i$. Le
groupe de décomposition $N_{s'}$ s'identifie à $G_{K^h_s}$, groupe de Galois du
corps des fractions de l'hensélisé (non strict) — le $\pi_1$ d'une « rondelle
trouée », ou mieux d'un fibré en disques épointés autour du germe de diviseur
$\overline{\{s\}}$ — et l'on a le diagramme de suites exactes, où $k = k(s)$ :

\begin{tikzcd}
 & & I \arrow[r, "\sim"] \arrow[d] & I_{s'} \arrow[d] & \\
1 \arrow[r] & \Delta \arrow[r] \arrow[d, no head, "="'] & \pi_1(U_{K^h_s}) \arrow[r] \arrow[d] & G_{K^h_s} \arrow[r] \arrow[d] & 1 \\
1 \arrow[r] & \Delta \arrow[r] & \pi_1(U_{\mathcal{O}^h_s}) \arrow[r] & G_{k} \arrow[r] & 1
\end{tikzcd}

\footnote{Le diagramme (26) de la page 87 a de plus des $1$ en haut et en bas
des colonnes ; nous les omettons. À la page 86, il écrit $\mathcal{E}_U$ là où
le diagramme écrit $\mathcal{E}_{U_{\tilde{\mathcal{O}}}}$.}La dernière ligne
s'identifie à la suite $1 \to \Delta \to \pi_1(U_s) \to G_{k(s)} \to 1$ de la
fibre $U_s$. Le noyau $I$ de la colonne médiane s'envoie isomorphiquement sur
$I_{s'} \simeq \widehat{\mathbb{Z}}(1)$. La section définie par $f$ donne un
relèvement
\[
\lambda : G_{K^h_s} \longrightarrow \pi_1(U_{\mathcal{O}^h_s}) \simeq
\pi_1(U_s)
\]
de $G_{K^h_s} \to G_{k(s)}$ ; et il « résume » :

\emph{Les conditions suivantes sont équivalentes : a) $\lambda$ est injectif ;
b) $\lambda|_{I_{s'}} = k_i \circ (\mu \cdot \mathrm{id})$ pour un
homomorphisme de lacets $k_i : \widehat{\mathbb{Z}}(1) \to \Delta$ et un
$\mu \in \mathbb{N}^{*}$ ; c) $\lambda|_{I}$ n'est pas trivial ; d) la section
de $\pi_1(U_{K^h_s})$ sur $G_{K^h_s}$ ne provient pas d'une section de
$\pi_1(U_{\mathcal{O}^h_s})$ ; e) $f$ n'est pas définie en $s$ ; f) $f(s) \in
X_s \setminus U_s$. Et alors l'indice $i$ de b) est $f(s)$.}

Ces équivalences reposent sur le critère local des pages 72 à 78 ; l'entier
$\mu$ est la multiplicité d'intersection $n$ de la page 77, c'est-à-dire la
multiplicité en $s$ du diviseur $f^{*}T$ (page 93). Dans ce cas le groupe de
décomposition $N_{s'} \simeq G_{K^h_s}$, extension de $G_{k(s)}$ par
$\widehat{\mathbb{Z}}(1)$, s'identifie à un sous-groupe d'indice $\mu$ du
normalisateur de $L_i$ dans $\pi_1(U_s)$\footnote{La page écrit
« $\mathcal{E}_{X_s}$ » pour le groupe ambiant, et « ext. de $\mathcal{E}_{k(s)}$
par $\pi$ » là où il faut « par $T$ », comme le dit la page 90.}. La section,
qui ne provient d'aucune section de $\pi_1(U_{\mathcal{O}^h_s})$, définit donc
« presque » un relèvement — seul le facteur $I_{s'}$ y fait obstacle — et toute
section de cette extension de $G_{k(s)}$ par $\widehat{\mathbb{Z}}(1)$ donne une
section de seconde espèce de $\pi_1(U_s)$ sur $G_{k(s)}$.

Plus généralement (pages 90 à 92), la section $f_K$ donne une famille de
sections des $\pi_1(U_s)$ sur les $G_{k(s)}$, pour $s$ dans le domaine de
définition $S'$ de $f$, par le même mécanisme : pour tout $s \in S$ et $s'$
au-dessus dans la normalisée de $S$ dans $\overline{K}$, on a $1 \to I_{s'} \to
N_{s'} \to G_{k(s)} \to 1$, avec $N_{s'} \simeq G_{K^h_s}$ et $I_{s'} \simeq
G_{K^{sh}_s}$\footnote{Il note $S'$ à la fois le domaine de définition de $f$
(page 90) et la normalisée de $S$ dans $\overline{K}$ (page 91), et $\tilde K$
les deux corps $K^h_s$ et $K^{sh}_s$.} ; si la composée s'annule sur
$I_{s'}$ — c'est dire $s \in S'$ —, on obtient $G_{k(s)} \to \pi_1(U_s)$, une
section. Pour $s \notin S'$, $D = S \setminus S'$ est purement de codimension
$1$ (image inverse de $T$ par $f$) ; au point générique $s_\alpha$ d'une
composante $D_\alpha$, $f$ donne un point rationnel $i_\alpha \in X_{s_\alpha}
\setminus U_{s_\alpha}$, et les sections de seconde espèce relatives à
$L_{i_\alpha}$, notamment celles que définit une équation locale de $D_\alpha$.
On voit comment elles se spécialisent aux points de $D_\alpha$ où la section
correspondante du fibré conormal reste régulière et inversible ; les autres
points demanderaient une étude analogue, « d'un intérêt mineur ».

Enfin (pages 95 à 98) : une section de première espèce $\varphi : G_K \to
\pi_1(U_K)$ définit, pour toute section de $G_K$ au-dessus d'un ouvert de
$G_{\mathbb{Q}}$ — ou de $G_{k'}$ pour un sous-corps $k' \subset K$ —, une
section correspondante de $\pi_1(U_K)$, d'où des actions extérieures sur
$\Delta$. Il veut les comparer à la famille des sections des $\pi_1(U_s)$ sur
les $G_{k(s)}$, où les $k(s)$ sont des corps résiduels et non des sous-corps.
Si $S$ est une courbe sur $k$, tout point fermé $s$ définit des groupes de
décomposition $N_{s'} \subset G_K$ au-dessus d'ouverts $G_{k(s)} \subset G_k$,
et l'on compare les sections de ces groupes à la section donnée : si $f$ est
définie en $s$, la section restreinte est de première espèce, induite par le
point $f(s) \in U_s(k(s))$ ; sinon elle est de seconde espèce et normalise un
$L_i$\footnote{La page 97 est d'écriture rapide ; plusieurs liaisons entre les
formules ne sont pas lues.}.

\pagerange{98}{103}
\subsection*{Mauvaise réduction (pages 98 à 103)}

Tout cela supposait bonne réduction. Hors de cette hypothèse — exemple type :
$K$ corps des fonctions d'une courbe modulaire pour les courbes elliptiques,
$s$ une pointe, où la courbe elliptique dégénère —, on passe à un trait
hensélien $S$ de corps résiduel $k$ (de type fini sur $\mathbb{Q}$), une courbe
$U_\eta$ de type $(g, \nu)$ (ici $(1,1)$) sur le point générique, et
\[
1 \to \widehat{\mathbb{Z}}(1) \to G_\eta \to G_k \to 1, \qquad 1 \to \Delta \to
\pi_1(U_\eta) \to G_\eta \to 1 .
\]
L'action extérieure de l'inertie sur $\Delta$ n'est pas triviale ; mais tout
scindage de la première suite définit une action extérieure de $G_k$ sur
$\Delta$. Est-elle « admissible », c'est-à-dire définie par une courbe
algébrique ? Sans doute pas — quelle courbe serait-ce, « la courbe réduite de
type $(g, \nu)$ (inexistante !) » ? En marge : « mais si !? ». Et il
note : « à tâter pour la structure à l'infini des variétés modulaires de
Teichmüller ! »\footnote{Lecture en partie lacunaire ; c'est le thème du
dossier 147, « Structure à l'infini des $M_{g,\nu}$ ».}. S'il faut hasarder une
conjecture dans « cette situation scabreuse » :

\textbf{Conjecture de bonne réduction.} \emph{Équivalence de : a) $U_\eta$ a
bonne réduction sur $S$ ; b) l'action extérieure de l'inertie sur $\Delta$ est
triviale ; c) l'action de l'inertie sur $\Delta^{\mathrm{ab}} =
H_1(U_{\bar\eta})$ est triviale ; d) de même sur $H_1(U_{\bar\eta},
\mathbb{Z}_\ell)$, $\ell$ premier à la caractéristique résiduelle ; e) $k$ étant
de type fini sur $\mathbb{Q}$, l'extension de $G_k$ par $\Delta$ définie par un
scindage est admissible (quitte à passer à un ouvert).}\footnote{Un alinéa f),
et une parenthèse en tête de la page 101, sont trop lacunaires pour être
rendus.}

On a a) $\Rightarrow$ b) $\Rightarrow$ c) $\Rightarrow$ d), et ces
conditions impliquent e) ; d) $\Rightarrow$ a) « devrait être plus ou moins
connu »\footnote{Ce ne l'est pas, et c'est faux en général : en genre $0$,
l'inertie agit trivialement sur $H_1 = \widehat{\mathbb{Z}}(1)^I /
\widehat{\mathbb{Z}}(1)$ même quand des pointes se rencontrent dans la fibre
spéciale ; en genre $\ge 2$, une courbe à réduction stable de type compact a une
jacobienne à bonne réduction. Pour une courbe elliptique, d) $\Rightarrow$ a)
est le critère de Néron–Ogg–Šafarevič. Ce qui est vrai en général est a)
$\Leftrightarrow$ b) : c'est le critère de bonne réduction d'Oda (1990, cas
propre), étendu par Tamagawa (1997) au cas affine.} ; mais il n'a « pas d'idée
pour le moment » pour voir que e) implique la bonne réduction.

Il dispose donc d'un principe pour construire des actions extérieures de
$G_k$ sur des groupes à lacets qui, sans doute, ne sont pas géométriques, par
des courbes « se réduisant mal » ; mais il n'arrive pas par cette voie à
construire des actions effectives, ou associées à des actions déjà
géométriques, qui ne soient pas elles-mêmes de l'un des deux types de la page
64. Le reste de la page 103 est blanc.

\pagerange{104}{135}
\section*{§4. Sections d'extensions, et anneaux de valuation généraux (pages 104 à 135)}

\pagerange{104}{114}
\subsection*{Révision des notations (pages 104 à 114)}

Pour $X$ connexe, $\pi_1(X)$ est le groupe fondamental profini comme groupe
extérieur ; un revêtement universel $\widetilde{X}$ le précise en un groupe
profini, $\pi_1(X; \widetilde{X}) = \mathrm{Aut}_X(\widetilde{X})$ ; un point
géométrique $\xi$ — un point $x$ et une extension séparablement close $\Omega$
de $k(x)$ — en donne un autre, $\pi_1(X, \xi)$, qui ne dépend que de la clôture
séparable de $k(x)$ dans $\Omega$\footnote{Il note ces groupes
$\mathcal{E}_X$, $\mathcal{E}_X^{\tilde X}$, $\mathcal{E}_X^{\xi}$,
$\mathcal{E}_X^{\Omega}$. Plusieurs mots de la page 105 ne se lisent pas ; le
sens est celui de la formule (4).}. Plus généralement tout $Y$-schéma $Z$
$1$-connexe au-dessus de $Y$ joue le rôle de foncteur fibre, et $(X,
\widetilde{X}) \mapsto \pi_1(X; \widetilde{X})$ est un foncteur. Pour $X \to Y$
et $T \to X$ avec $T$ $1$-connexe on pose
\[
\pi_1(X/Y; T) = \mathrm{Ker}\bigl(\pi_1(X; T) \to \pi_1(Y; T)\bigr) .
\]
S'il y a une « suite exacte d'homotopie universelle », par exemple si $Y =
\operatorname{Spec} K$ et $X$ géométriquement connexe sur $K$, on a
\[
1 \longrightarrow \pi_1(X/Y; T) \longrightarrow \pi_1(X; T) \longrightarrow
\pi_1(Y; T) \longrightarrow 1, \qquad \pi_1(X/Y; T) \simeq \pi_1(X \times_Y T;
T),
\]
le groupe fondamental des fibres géométriques ; sinon $X \to Y$ se factorise en
$X \to Y' \to Y$, $Y'$ pro-étale fini, avec $\pi_1(X) \to \pi_1(Y')$
surjectif\footnote{La formule (16) porte $X/Y$ et non $X/Y'$, et la phrase qui
suit écrit $X \times_Y T$ là où l'on attend $X \times_{Y'} T$.}. Les groupes
fondamentaux des diverses fibres géométriques sont canoniquement isomorphes
entre eux ; mais $\pi_1(X/Y)$ n'est qu'une classe d'isomorphie de groupes
profinis, pas même un groupe extérieur, l'arbitraire venant de l'action de
$\pi_1(X)$ par automorphismes intérieurs. Ce point de vue regarde $\pi_1(X)$
comme une \emph{extension} de $\pi_1(Y)$, qui joue le rôle d'un groupe de
Galois, par $\pi_1(X/Y)$, qui joue le rôle d'un groupe fondamental
géométrique. Le cas principal est $Y = \operatorname{Spec} K$, où $\pi_1(Y) =
G_K$.

\pagerange{114}{120}
\subsection*{Modèles, inertie, groupe générique (pages 114 à 120)}

Soient $L$ de type fini sur un corps $K$ et $X$ un modèle propre et régulier de
$L/K$. Alors $\pi_1(X)$ est un quotient de $G_L$ qui ne dépend pas du modèle
choisi ; c'est la partie « universellement générique » de $G_L$, qui classe les
extensions finies de $L$ non ramifiées sur tout modèle régulier, propre ou
non\footnote{C'est-à-dire non ramifiées en toute valuation divisorielle. Pour
la caractéristique $0$, invariance birationnelle de $\pi_1$ des modèles propres
réguliers et pureté de Zariski–Nagata.}. Pour un modèle quelconque $U = X
\setminus Z$, on a des surjections $G_L \to \pi_1(U) \to \pi_1(X)$, et
$\pi_1(X)$ est le quotient de $\pi_1(U)$ par le sous-groupe distingué fermé
engendré par les inerties $I_{x'_i}$ aux points génériques $x_i$ des
composantes de codimension $1$ de $Z$. On a $G_L \simeq \varprojlim_U \pi_1(U)$
et $\Delta_L \simeq \varprojlim_U \Delta_U$.

Pour un anneau de valuation discrète $V$ de $L/K$ correspondant au point
générique d'un diviseur d'un modèle — une \emph{valuation divisorielle} —, de
point fermé $x$, on a le groupe de décomposition $N_{x'}$ et le groupe
d'inertie $I_{x'}$, avec
\[
N_{x'}/I_{x'} \simeq G_{k(x)}, \qquad I_{x'} \simeq \widehat{\mathbb{Z}}(1)
\subset \Delta_L ,
\]
$k(x)$ étant de degré de transcendance $n - 1$ sur $K$ ; et $N_{x'} \simeq
\pi_1(\operatorname{Spec} L^h_x)$ (hensélisé non strict), $I_{x'}$ étant le
noyau de $N_{x'} \to G_{k(x)}$. Un passage marqué d'un double trait et d'un
point d'interrogation note $N_{x'} = \mathrm{Norm}_{G_L}(I_{x'})$, « lié à
l'idée qu'on devrait avoir $I_{x'} = \mathrm{Norm}_{\Delta_L}(I_{x'})$ — à tirer
au clair à l'occasion ! »\footnote{La seconde égalité ne peut valoir qu'en
degré de transcendance $1$ : le normalisateur de $I_{x'}$ dans $\Delta_L$
contient $N_{x'} \cap \Delta_L$, extension par $I_{x'}$ de la partie
géométrique de $G_{k(x)}$, non triviale dès que $n \ge 2$. Pour $n = 1$, elle
est vraie.}.

Si l'homomorphisme $\pi_1(\operatorname{Spec} L^h_x) \to \pi_1(U)$ se
factorise par $\operatorname{Spec} \mathcal{O}^h_x$ — si $V$ a un centre sur
$U$ —, l'image de $I_{x'}$ dans $\pi_1(U)$ est nulle ; la réciproque est
posée en question. Avec un modèle propre $X \supset U$, la condition est
simplement que le centre de $V$ sur $X$ soit dans $U$. Ainsi $\pi_1(U)$ se
reconstruit à partir de $G_L$, comme quotient par le sous-groupe distingué
fermé engendré par les $I_{x'}$ des valuations divisorielles \emph{sans}
centre sur $U$\footnote{La page 120 dit « en prenant tous les $V$ ayant un
centre sur $U$ », c'est-à-dire les $V = \mathcal{O}_{U,x}$, et en divisant par
leurs inerties ; c'est l'inverse de ce qu'exige ce qui précède, où les
inerties tuées sont celles des points de $X \setminus U$. La phrase de la page
est d'ailleurs en partie incertaine ; nous donnons la forme cohérente avec la
page 119.}.

\pagerange{120}{123}
\subsection*{L'espace de Riemann–Zariski (pages 120 à 123)}

Il introduit l'espace
\[
\mathfrak{X}_{L/K} = \varprojlim_{X \text{ modèle propre de } L/K} X ,
\]
aujourd'hui l'\emph{espace de Riemann–Zariski} de $L/K$. Pour $x \in
\mathfrak{X}$, $k(x)$ est de degré de transcendance $d \le n$ sur $K$, sa
\emph{dimension} ; $n - d$ est sa codimension. Il y a un seul point de
codimension $0$, et les points de codimension $1$ correspondent exactement aux
valuations divisorielles de $L/K$. Sauf pour $n \le 1$, $\mathfrak{X}$ n'est pas
un schéma\footnote{La page ajoute, « sauf erreur », que ces deux types de points
sont les seuls dont l'anneau local soit noethérien. C'est faux pour $n \ge 2$ :
il existe des valuations discrètes de rang $1$ non divisorielles, par exemple
celle d'un arc transcendant ($v(F) = \mathrm{ord}_t\, F(t, \varphi(t))$ sur
$K(x, y)$, avec $\varphi$ série formelle transcendante), dont l'anneau est
noethérien et dont le corps résiduel est algébrique sur $K$.}.

\textbf{Proposition.} \emph{Pour tout $x \in \mathfrak{X}$, l'anneau local
$\mathcal{O}_{\mathfrak{X},x}$ est un anneau de valuation de corps des
fractions $L$ contenant $K$, et $x \mapsto \mathcal{O}_{\mathfrak{X},x}$ est une
bijection de $\mathfrak{X}$ sur l'ensemble de ces anneaux de valuation.}

C'est le théorème de Zariski. L'esquisse de la page : a) pour $f \notin
\mathcal{O}_{\mathfrak{X},x}$, écrire $f = g/h$ sur un modèle, et, par
éclatements, séparer les diviseurs de $g$ et $h$, pour obtenir $1/f \in
\mathcal{O}_{\mathfrak{X},x}$\footnote{La page invoque la résolution des
singularités ; elle n'est pas nécessaire, l'éclatement de l'idéal $(g, h)$
suffit. L'esquisse est en partie illisible.} ; b) un anneau de valuation $V$
domine un point bien déterminé de chaque modèle propre (critère valuatif), d'où
un $x$ tel que $V$ domine $\mathcal{O}_{\mathfrak{X},x}$, donc lui est égal.
Tout modèle $U$ définit un ouvert de $\mathfrak{X}$, mais les ouverts ainsi
obtenus ne sont pas stables par réunion finie.

\pagerange{123}{127}
\subsection*{Germes de sections, et une conjecture birationnelle (pages 123 à 127)}

Pour tout anneau local normal $\mathcal{O} \subset L$ de corps des fractions
$L$ on a encore des groupes de décomposition et d'inertie. L'idée : pour un
anneau de valuation contenant $K$, les groupes d'inertie ont une structure
simple, « des $\widehat{\mathbb{Z}}(1)^{\alpha}$ essentiellement »\footnote{La
page dit « où $\alpha$ est la codimension » du point. En caractéristique
résiduelle $0$, l'inertie est $\mathrm{Hom}(\Gamma_v, \widehat{\mathbb{Z}}(1))$
où $\Gamma_v$ est le groupe des valeurs : de rang $\alpha$ si $\Gamma_v \simeq
\mathbb{Z}^\alpha$, et $\alpha$ est le rang rationnel, majoré par la
codimension (inégalité d'Abhyankar), égal pour les valuations dites
d'Abhyankar seulement.} ; et si $x$ est de dimension $0$ ($k(x)$ fini sur $K$),
le groupe de décomposition est une extension de $G_{k(x)}$, ouvert dans $G_K$,
par l'inertie, qu'un choix de paramètres scinde. D'où un principe général
pour construire des \emph{germes de sections} de $G_L$ sur $G_K$ — sections
au-dessus d'un sous-groupe ouvert — à partir des points de dimension $0$ de la
« variété de Riemann » $\mathfrak{X}$.

Pour une sous-extension $L_1$ de $L/K$ de degré de transcendance $1$, on en
déduit des germes de sections de $G_{L_1}$ sur $G_K$\footnote{La page écrit
« sous-extension $L_1$ de $K$ » ; il s'agit d'une sous-extension de $L/K$.} ; en
trouve-t-on plus que celles des places de $L_1/K$ ? Non : $V_1 = V \cap L_1$
est un anneau de valuation de $L_1$, de corps résiduel fini sur $K$ (et $V_1 \ne
L_1$, puisque le corps résiduel de $V$ est algébrique sur $K$), donc une place ;
$G_L \to G_{L_1}$ envoie $N_V$ dans $N_{V_1}$ et $I_V$ dans $I_{V_1}$ ; le germe
obtenu normalise $I_{V_1}$, et sa restriction à $G_K^{\circ}$ le centralise,
puisque $G_K$ agit sur $I_{V_1} \simeq \widehat{\mathbb{Z}}(1)$ par $\chi$.

\textbf{Conjecture.} \emph{Soient $K \subset L$ de type fini sur $\mathbb{Q}$.
a) Toute section (ou germe de section) de $G_L$ sur $G_K$ normalise un $I_V$
associé à un anneau de valuation $V$ de $L/K$ à corps résiduel algébrique sur
$K$, uniquement déterminé ; elle provient d'une section du groupe de
décomposition $N_V$. b) Pour tout modèle élémentaire $U$ de $L/K$, toute
section de $\pi_1(U)$ sur $G_K$ se relève en une section de $G_L$ sur $G_K$.}

Une section de $\pi_1(U)$ se décrit donc par un anneau de valuation $V$ ; s'il
a un centre $x \in U(K)$, la section est celle du point $x$, et l'on conjecture
encore
\[
\Delta_U^{\,G_K^{\circ}} = 1,
\]
qui équivaut à la même relation pour $\pi_1(U)$ dès que le centralisateur de
$G_K^{\circ}$ dans $G_K$ est trivial, et qui se déduirait par dévissage de la
conjecture analogue en dimension relative $1$ (page 48)\footnote{La page écrit
ces deux relations, (32) et (33), pour $G_L$ et $\Delta_{L/K}$. Lues ainsi, elles
contredisent la page 125 et la conséquence conjecturale de la page 44 : la
section normalise $I_V \ne 1$, que $G_K^{\circ}$ centralise. Le contexte — une
section de $\pi_1(U)$ définie par un point $x \in U$, et le dévissage vers la
dimension relative $1$ — impose de les lire pour le modèle $U$, comme à la page
48. Une longue note en marge, en deux blocs, ne se lit que par fragments.}.

\pagerange{127}{131}
\subsection*{L'inertie le long du bord (pages 127 à 131)}

Pour comprendre b), il faut voir ce que deviennent dans $\pi_1(U)$ les
sections de $G_L$ associées à un $V$, et d'abord le sous-groupe $\Delta_V
\subset \Delta_U$ image de $I_V$, dont il « devine » qu'il donne une réalisation
algébrique d'un groupe de lacets. Des $V$ distincts donnent-ils des $\Delta_V$
distincts ?

Soit $X \supset U$ un modèle propre tel que $D = X \setminus U$ soit un
diviseur à croisements normaux de composantes $D_i$ régulières, $i \in I$. Pour
$x \in D$, soient $I(x) = \{i \mid x \in D_i\}$, $\mathcal{O}^h_x$ l'hensélisé
de $\mathcal{O}_{X,x}$ et $U^h_x$ l'image inverse de $U$. En caractéristique
$0$, par le lemme d'Abhyankar,
\[
1 \longrightarrow \widehat{\mathbb{Z}}(1)^{I(x)} \longrightarrow
\pi_1(U^h_x) \longrightarrow G_{k(x)} \longrightarrow 1 .
\]
Si le centre de $V$ sur $X$ est $x \in D$, l'homomorphisme $N_V \to \pi_1(U)$ se
factorise en $N_V \to \pi_1(U^h_x) \to \pi_1(U)$, l'image du second étant le
groupe de décomposition de $x$. Sous réserve de vérifications, il présume :

\begin{enumerate}
\item[1°)] que $\pi_1(U^h_x) \to \pi_1(U)$ est injectif, c'est-à-dire que
$\widehat{\mathbb{Z}}(1)^{I(x)} \to \Delta_U$ l'est ;
\item[2°)] que si $V$ est de dimension $0$, $N_V \to \pi_1(U^h_x)$ a une image
ouverte, c'est-à-dire que $I_V \to \widehat{\mathbb{Z}}(1)^{I(x)}$ est
surjectif.
\end{enumerate}

Une note en marge, d'une encre plus foncée, ajoutée après coup, dit de 1°) :
« C'est sûrement “en général” \emph{faux} pour $\operatorname{card} I(x) \ne
1$, comme on voit en remplaçant $X$ par un $X'$ dominant $X$, en éclatant […]
un point de $D$… \textbf{Mais} est-il vrai qu'on [peut trouver] $X \supset U$
tel que [ce soit] localement injectif ? »\footnote{Mots entre crochets : lecture
incertaine de la transcription, reconstruits par nous pour la phrase. La page
133 en fournit la raison : voir la note sur $\pi_2(D_J)$.} De 2°), il dit qu'elle
est purement locale et « doit être essentiellement triviale, que ma
présomption s'avère vraie ou fausse » ; elle est fausse en
général\footnote{Contre-exemple : en dimension $2$, avec $I(x) = \{1, 2\}$ et
des équations locales $t_1$, $t_2$, la valuation d'arc $v(F) =
\mathrm{ord}_t\, F(t, t\varphi(t))$, $\varphi$ transcendante avec $\varphi(0)
\ne 0$, est de dimension $0$ et vérifie $v(t_1) = v(t_2) = 1$ ; $I_V \simeq
\widehat{\mathbb{Z}}(1)$ s'envoie diagonalement dans
$\widehat{\mathbb{Z}}(1)^2$, d'image non ouverte. La surjectivité vaut quand
$v(t_i)_{i \in I(x)}$ engendrent un facteur direct du groupe des valeurs,
par exemple pour la valuation monomiale de poids $1$ et $\sqrt{2}$.}. 1°) au contraire est « une question
de nature globale sur $U$, et sans doute aucunement triviale ».

Si ces énoncés valaient, les sections de $\pi_1(U)$ associées à une valuation
de corps résiduel $K$ normaliseraient un sous-groupe $\Delta_x \simeq
\widehat{\mathbb{Z}}(1)^{I(x)}$ de $\Delta_U$, non trivial si la valuation n'a
pas de centre sur $U$ — « ce qui est justement dans le sens des conjectures
(qui pouvaient d'abord sembler hérétiques !) du §2 ».

\pagerange{131}{135}
\subsection*{Les strates du bord (pages 131 à 135)}

Avant d'aborder 1°) et 2°), il regarde les classes de conjugaison des
sous-groupes $I_x$ pour $x \in X \setminus U$ variable. Pour $J \subset I$,
soient
\[
D_J = \bigcap_{j \in J} D_j, \qquad D_J^{*} = \{x \in X \mid I(x) = J\} =
D_J \setminus \bigcup_{i \notin J} D_i ;
\]
$D_J$ est régulier, de codimension $\operatorname{card} J$, peut-être vide.
Les $J$ tels que $D_J \ne \emptyset$ forment un ensemble simplicial sur $I$ —
c'est aujourd'hui le \emph{complexe dual} du diviseur $D$. « Quelle est sa
signification topologique ? » Ses composantes connexes sont celles de $D$, et
celles du « bout » $\widehat{X} \setminus D$ ($\widehat{X}$ le complété formel
de $X$ le long de $D$), notion qui ne dépend pas de la compactification ; si les
$D_J$ étaient « $\infty$-connexes », l'ensemble simplicial exprimerait le type
d'homotopie de $U$ à l'infini\footnote{Le type d'homotopie du complexe dual
ne dépend en effet que de $U$ (Danilov ; Thuillier, 2007, par l'espace de
Berkovich ; Payne).}.

Soit $X_J$ un « localisé » de $X$ le long de $D_J$ (en marge : un voisinage
« tubulaire » ou « formel » vaudrait mieux), et $U_J = X_J \setminus D_J$. On
trouve une suite
\[
\widehat{\mathbb{Z}}(1)^J \longrightarrow \pi_1(U_J) \longrightarrow
\pi_1(D_J) \longrightarrow 1,
\]
début de la suite d'homotopie d'un fibré en produits de disques épointés
au-dessus de $D_J$. Il a « eu un peu peur d'y mettre zéro à gauche, à cause du
$\pi_2(\overline{D}_J) \to \pi_1(\text{fibre})$ (?) ! » En comparant avec la
suite $1 \to \Delta_U \to \pi_1(U) \to G_K \to 1$, il voit que la composée
$\pi_2(D_J) \to \widehat{\mathbb{Z}}(1)^J \to \Delta_U$ est nulle : si
$\widehat{\mathbb{Z}}(1)^J \to \Delta_U$ est injectif (1°), $\pi_2(D_J) \to
\widehat{\mathbb{Z}}(1)^J$ est nul\footnote{Sa peur est fondée. En topologie,
ce bord est donné par les classes de Chern des fibrés normaux des $D_j$ le long
de $D_J$, et il n'est pas nul en général : pour une courbe exceptionnelle
$D \simeq \mathbb{P}^1$ d'auto-intersection $-1$ sur une surface, le voisinage
épointé a le type d'homotopie de $S^3$, et le lacet autour de $D$ est trivial.
C'est la même situation qui rend 1°) faux après éclatement, comme le dit la
marge de la page 129.}. Il admet donc l'exactitude à gauche, et obtient un
homomorphisme de suites exactes

\begin{tikzcd}
1 \arrow[r] & \widehat{\mathbb{Z}}(1)^J \arrow[r] \arrow[d] & \pi_1(U_{J,\alpha}) \arrow[r] \arrow[d] & \pi_1(D_{J,\alpha}) \arrow[r] \arrow[d] & 1 \\
1 \arrow[r] & \Delta_U \arrow[r] & \pi_1(U) \arrow[r] & G_K \arrow[r] & 1
\end{tikzcd}

pour chaque composante connexe $D_{J,\alpha}$ de $D_J$\footnote{La page porte
deux numéros devant ce diagramme, « (38') » surchargé et « (38) », et
emploie (39) deux fois, aux pages 134 et 135.}. Les homomorphismes
$\widehat{\mathbb{Z}}(1)^J \to \Delta_U$ induits par les divers $x \in
D_{J,\alpha}$ sont tous conjugués entre eux, et l'on peut faire les choix de
façon cohérente, de sorte que l'indétermination ne porte plus que sur l'image
de $\pi_1(D_{J,\alpha}) \to G_K$. Pour $J = \{i\}$, $D_J = D_i$ est connexe, et
l'on obtient
\[
\kappa_i : \widehat{\mathbb{Z}}(1) \longrightarrow \Delta_U ,
\]
homomorphisme extérieur bien déterminé, compatible aux actions extérieures de
$G_K$ sur $\widehat{\mathbb{Z}}(1)$ et sur $\Delta_U$ — l'analogue, en dimension
quelconque, des homomorphismes de lacets $k_i$ des courbes.

Le dossier s'arrête sur un paragraphe barré de grandes croix, qui commençait
à comparer les composantes de $D_{J'}$ et de $D_J$ pour $J' \supset J$\footnote{Une
longue note en marge, sur toute la hauteur de la page 135, paraît préciser les
choix de revêtements universels qui rendent les $\kappa_i$ bien définis ; seuls
quelques mots se lisent.}.

\end{document}
